\chapter[数据结构化札记]{数据结构化札记\footnote{本文以 1970 年在 Marktoberdorf 举办的北约暑期学校上的一系列讲座为基础。}}

{
	\noindent
	\scshape\hfill\scriptsize C. A. R. HOARE\hfill
}
\renewcommand{\leftmark}{\normalfont\scriptsize\hfill C. A. R. HOARE\hfill}
\medskip

\section{引言}

在认识复杂现象的过程中，人类心智能够运用的最强大工具就是抽象。抽象始于这样的认识：现实世界中的某些对象、情境或过程存在相似之处；于是我们决定把注意力集中在这些相似之处，暂时忽略它们之间的差异。一旦发现哪些相似性与预测和控制未来事件有关，我们就会倾向于把这些相似性看作根本，把差异看作次要。这时可以说，我们已经形成了一个涵盖这些对象或情境的抽象概念。到了这个阶段，通常会引入一个词语或图形作为这一抽象概念的符号；每一次说出或写下这个词语或图形，都可以用来\textit{表示}相应情境的某个特例或一般情形。

表示的主要用途，是把现实世界中重要方面的信息传达给别人，并把这些信息记录下来：一方面帮助记忆，另一方面传给后代。不过，在原始社会中，人们有时相信表示本身就有用，因为他们以为，操纵表示本身就能使现实世界发生相应的变化。因此，我们听说过这样的做法：把针扎进敌人的蜡像，以使真人的相应部位疼痛。这类活动是魔法和巫术的典型特征。现代科学家则认为，操纵表示可以用来预测现实世界中的事件及变化的结果，却不能引起这些事件和变化。例如，通过操纵某些函数和方程的符号表示，他可以预测下落物体撞击地面时的速度，同时也知道，这既不会使物体下落，也不会在物体确实落下时减轻最后的撞击。

抽象过程的最后一个阶段要精细得多：尝试用简洁而有力的\textit{公理}，概括某一抽象所涵盖的情境和对象的最一般事实；并在这些公理正确描述现实世界的条件下，严格证明操纵表示所得到的结果也能成功地应用于现实世界。例如，欧几里得几何的公理与真实、可测量的世界足够接近，因而几何作图和定理可以合理地用于土地测量以及地表勘测等实际工作。

因此，抽象过程可以概括为四个阶段：

\begin{enumerate}[wide, nosep, label=(\arabic*)]
	\item 抽象：决定关注现实世界中许多对象或情境共有的性质，忽略它们之间的差异。
	\item 表示：选择一组符号来代表这一抽象；它可以作为交流的手段。
	\item 操纵：规定符号表示的变换规则，以预测对现实世界进行类似操纵所产生的效果。
	\item 公理化：严格陈述从现实世界中抽象出来的性质；这些性质既为现实世界中的操纵所共有，也为对其表示符号的操纵所共有。
\end{enumerate}

\subsection{数与数码}
让我们用一个比较具体的例子\textemdash{}数四\textemdash{}来说明上述颇为抽象的描述。在现实世界中，人们注意到，对象可以聚在一起构成集合，例如四个苹果。这已经需要某种抽象活动，也就是决定暂时忽略集合中各个苹果之间的差异\textemdash{}例如，其中一个已经坏了，两个还没熟，第四个已被鸟啄掉了一部分。

现在可以考虑几个不同的集合，每个集合都有四个对象，例如四个橘子、四个梨、四根香蕉等等。如果选择忽略这些集合之间的差异，而把注意力集中在它们的相似之处，就可以形成数四这一比较抽象的概念。同样的过程可以产生数 3、15 等概念；再进一步抽象，就会形成自然数这一概念。

接下来考虑如何表示这个概念，例如写在羊皮纸上，或刻在石头上。数的表示称为数码（numeral）。早期的罗马数码显然具有图画性质，就是刻在石头上的四道竖线：IIII。另一种更方便的表示是 IV。阿拉伯数码（十进制）的图画性质较弱，但同样有选择余地：4 和 04（当然还有 004 等等）都被认为是有效的数码，表示同一个数。

再来看一种极便于处理的表示，前提是处理者是一台电子数字计算机。在这里，数四用一组铁氧体磁芯的不同磁化方向表示。这些磁化状态有时又用行式打印机纸上的 0、1 序列表示，也就是这个数的二进制表示。

操纵数码的一个简单例子是加法，它可以用来预测现实世界中把两组对象合在一起的结果。罗马数码的加法规则非常简单、直观，也容易应用。十以内的阿拉伯数码加法规则相当不直观，必须学习；但对于远大于十的数，它比罗马数码的方法更方便。二进制表示的加法并不适合人来做；但对计算机而言，这却是最简单、最好的表示。因此，在多种表示之间作选择时，可以考察每一种表示在特定环境中是否易于操纵。

最后来到公理化阶段。最广为人知的自然数公理系统是皮亚诺公理，它在上世纪末才首次提出，而那时自然数早已被普遍使用。在当今时代，抽象数学思想的公理化通常更紧密地跟随这些思想的发展；实际上，它还可以帮助澄清概念，防止混淆和错误，并向别人说明概念的本质特征。如果把程序所依据的前提和公理严格表述出来，或许就能减少当今程序设计实践中常见的混淆和错误，也能帮助编写文档，向别人解释程序及程序设计概念。

\subsection{抽象与计算机程序设计}

我相信，把数学应用于现实世界所依赖的抽象过程，也正是把计算机应用于现实世界所依赖的过程。设计程序的首要要求，是关注情境中有关的特征，忽略那些被认为无关的因素。例如，分析飞机某种拟议机翼设计的颤振特性时，被认为有关的是它的弹性；颜色、形状和制造工艺则被认为无关，除非这些因素影响了弹性。再举一个商业方面的例子：公司的员工有许多身体和心理特征，而在为公司设计工资核算程序时，这些特征会被忽略。

程序设计的下一个阶段，是决定如何在计算机中表示抽象出来的信息。弹性函数可以用适当数量的离散点上的函数值表示，而这些值又可以通过多种方式表示为一个二维数组。另一种办法是用计算得到的函数给出弹性，数据则保存为这个函数的多项式系数或切比雪夫系数向量。计算机中的工资文件包含若干记录，每条记录对应一名员工。对于记录中的每个有关属性，都必须在程序设计中决定其表示方式。

公理化通常不被视为程序设计中的独立阶段，而且往往只是隐含存在。对于飞机颤振，公理化就是建立一组微分方程：假定它们描述了实际机翼对某些应力的反应，并希望它们也能描述计算机上求近似解的过程。对于工资核算，公理对应的是对现实世界各个方面的描述，而这些描述需要体现在程序中\textemdash{}例如，实发工资等于应发工资减去各项扣款。

最后的任务是编写程序，让计算机对数据表示进行操纵，使这些操纵对应于我们所关心的现实世界中的操纵。

程序的成功取决于三个基本条件：

\begin{enumerate}[wide, nosep, label=(\arabic*)]
	\item 公理化正确地描述了它所关注的现实世界的有关方面。
	\item 公理化正确地描述了程序的行为，也就是说，程序没有错误。
	\item 表示方式和操纵方法的选择，使程序在计算机上的运行成本可以接受。
\end{enumerate}

为了简化计算机程序的设计和开发任务，最好能把这三个阶段适当地分开，并按合适的顺序进行。第一阶段（公理化）的成果，是用严格的逻辑语言陈述关于现实世界的前提，并表述程序需要实现的目标。第二阶段的成果，是一个算法或抽象程序，可以证明它在给定前提下能够完成所陈述的任务。第三阶段则要决定，如何在计算机的存储器中表示和操纵各种数据项，以达到可以接受的效率。只有这三个阶段都圆满结束以后，才开始最后的编码和测试阶段，把选定的算法落实为对选定数据表示进行操作的程序。

当然，这是一幅略显理想化的图景：程序设计这一智力工作，从问题的抽象表述稳步推进到求解中越来越具体的方面。在实践中，即使是在表述问题时，程序员也必须对问题能否解决有某种直觉；设计抽象程序时，也必须觉得存在一种效率足够高的表示。这些直觉和感觉常常会出错，
而对表示的深入研究，甚至最后的编码，都会迫使我们退回到较早阶段，有时甚至需要彻底改变解决问题的思路。不过，凭直觉预先思考并承担失败的风险，是一切具有发明和构造性质的智力活动的特征。这并不损害有条理地起步、把阶段分得相对清楚所带来的好处。

这种推进方式最重要的特点之一，是把实际编码推迟到几乎确信设计的其他方面都已成功完成之后。编码和程序测试通常是程序开发中成本最高的阶段，因此不希望在这一阶段开始以后再作修改。最好事先确保在最后阶段不会再有其他问题，例如，程序解决的是正确的问题，算法是正确的，程序各部分能协同完成总体任务，数据表示也具有足够的效率。本文的目的，就是探索获得这种把握的方法。

\subsection{高级程序设计语言中的抽象}

采用合适的高级程序设计语言，可以进一步发挥抽象在计算机程序设计和开发中的作用。事实上，使用高级语言而不是机器码所带来的好处，很大程度上可能正是因为这些语言纳入了成功的抽象，尤其是数据抽象。在计算机硬件和机器码程序员看来，每个数据项都只是一组位。但对于 ALGOL 60 或 FORTRAN 程序员来说，数据项是整数、实数、向量或矩阵；这些抽象也正是这两种语言最初面向的数值应用领域的基础。当然，语言实现者已把这些抽象概念映射为特定计算机上的具体位模式表示。但程序员在设计算法时，不必考虑这些对于他的目的大多无关的细节，工作也因此大为简化。

高级程序设计语言的另一个主要优点，即机器无关性，同样来自其抽象的成功。抽象不仅可以表达不同现实情境的重要特征，也可以表达这些情境在不同计算机表示中的重要特征。这样，每位实现者都可以选择一种表示，使其在自己的计算机上得到最高的操纵效率。

高级语言的第三个主要优点，是显著减少了出现程序设计错误的机会。用机器码编程时，很容易犯下愚蠢的错误，例如对浮点数使用定点加法，对布尔标志进行算术运算，或者让修改后的地址越界。采用高级语言时，可以通过三种方式防止这类错误：

\begin{enumerate}[wide, nosep, label=(\arabic*)]
	\item 使用错误算术指令这类错误，在逻辑上不可能发生。例如，用 ALGOL 表达的任何程序都不可能导致生成这类错误代码。
	\item 对布尔标志进行算术运算这类错误，会被编译器立即发现，绝不会在可执行程序中造成问题。
	\item 下标越界这类错误，可以在运行时检查数组下标的范围而发现。
\end{enumerate}

运行时检查虽常常必不可少，但几乎不可避免地比前两类检查成本更高、使用更不方便。因此，高级语言的设计应当扩大这两类错误的范围：在逻辑上不可能发生的错误，以及一旦发生就能由编译器发现的错误。事实上，巧妙的语言设计可以让大多数下标检查不损害运行时效率。

程序设计错误的自动预防和检测，同样可以归功于对抽象的成功运用。高级程序设计语言允许程序员声明自己对所用变量的值的类型有何意图，由此规定适用于该类型的值的操作的含义。这样，编译器就比较容易检查程序的一致性，防止错误进入执行阶段。

\subsection{记法}

介绍数据结构化理论，需要引入一些方便的记法，以表达其中的抽象。这些记法在很大程度上以数学家、逻辑学家和程序员已经熟悉的记法为基础。它们也经过设计，可以直接表达计算机算法，并尽量减少运行中的程序发生编程错误的机会。最后，这些记法的设计还确保了在数字计算机上存在高效的数据表示。既然它们的用途之一是表达算法，人们自然可能认为，它们构成了某种程序设计语言，而且应当编写一个自动翻译程序，把用这种语言表达的程序转换为计算机的机器码，从而消除程序开发中昂贵又易出错的编码阶段。

但这样的结论完全误解了引入这些记法的原因，而且可能产生一些很不理想的后果。最糟糕的后果，是人们可能以这种语言没有可用编译器、即使有了也难以广泛接受为由，拒绝本文阐述的程序设计方法所带来的主要好处。

不应把这些记法视为程序设计语言，是有充分理由的。某些操作（例如序列串接）虽然对抽象程序的设计及其性质的描述很有帮助，但用于计算机上的大型数据对象时，效率却低得惊人。在从抽象程序过渡到具体程序时，消除这类操作是程序设计过程中不可缺少的一部分。消除它们有时需要对算法和表示都作相当彻底的改变，一般不能交给自动翻译程序。如果这类代价高昂的运算符被纳入一种准备自动编译的语言，许多程序员很可能意识不到，自己有责任在开始使用计算机之前把它们消除。即使想这样做，他们也很难判断哪些替代表示和操作更为经济。在作这些至关重要的决定时，程序设计语言与最终机器比较接近，实际上是有帮助的：也就是说，能够根据相应源语言代码的形式和长度，直接预测机器码的效率。

自动实现这些记法还存在一种更隐蔽的危险：优秀的程序员很快就会知道，其中一些记法的效率明显低于其他记法，于是即使在抽象程序中也避免使用它们。这会形成某种思维障碍，可能严重影响他的创造能力。同样严重的是，一组固定记法的实现，也可能妨碍使用者根据自己对特定问题的理解，引入所需的记法和概念。

因此，必须明确区分两种语言：一种算法语言用于帮助定义、设计、开发程序并编写文档；另一种程序设计语言则用于最终把程序交给计算机。本文只讨论前一种语言。所有示例算法都用这种语言表达；这些程序的实际编码留作读者的练习。读者可以为此选择任何熟悉的语言：ALGOL、FORTRAN、COBOL、PL/I、汇编语言，或它们的任何可用组合。要理解各种数据表示各自的优点，就必须认识到它们会给最终代码带来什么影响。

虽然我如此明确地声明，自己并不打算再设计一种程序设计语言，但也必须承认：如果用来为算法编码的程序设计语言，实际上是设计该算法所用语言的一个\textit{子集}，就会带来一些好处。我还必须承认，所提出的算法语言中，存在一个相当大的子集，可以在当今普通计算机上以极高的效率实现，无论编译时还是运行时都是如此；而设计能够高效实现更大子集的计算机这一挑战，也可能在将来被接受。但即使没有这种子集的实现，也丝毫不影响把整套记法用作抽象程序设计工具的好处。

\subsection{小结}

本引言概述了下文的动机和总体思路。和通常一样，在读完本文其余部分以后再读一遍引言，会有更大的收获。

第二节解释类型的概念。这一概念对数据结构化理论至关重要；本节还把它与计算机程序设计实践中的有关操作和表示联系起来。

之后各节讨论具体的数据结构化方法，从简单结构逐步讲到更复杂的结构。

每种结构都会借助例子作非形式化的说明。随后，列出可以合法地用于这一结构的一组基本操作，以此定义对它的操纵。最后，给出若干可能的计算机表示，并说明每次选择合适表示时应当考虑的准则。

第 \ref{sec:example-examination} 节专门讨论一个例子，即编排考试时间表的程序。最后一节通过建立表达数据结构基本性质的公理，把整个论述置于严格的理论基础之上。这一节可以作为理论小结，用于查阅以深化理解，也可以作为证明程序正确性的基础。
\section[类型的概念]{类型的概念}

这里提出的数据结构化理论，在很大程度上依赖类型的概念。数学家、逻辑学家和程序员都熟悉这一概念。

(1) 在数学推理中，通常会相当明确地区分个体、个体的集合、集合族等等；也会区分实函数、复函数、泛函、函数的集合等等。实际上，数学家在推理中每引入一个新变量，通常都会立即说明它可以代表哪一类对象，例如：

\quad “设 $f$ 是两个实变量的实函数。”

\quad “设 $S$ 是整数集合组成的一个集合族。”

数学著作有时会给出一条总规则，把符号的类型与某种印刷字体联系起来，例如：

\hangindent=1.9\parindent \quad“我们用小写罗马字母表示个体，用大写字母表示个体的集合，用花体大写字母表示集合族。”

数学家一般不会用这类类型约定来作任意的区分；例如，通常不会用它来区分素数与非素数，或阿贝尔群与一般的群。实际上，数学家采用的类型约定，与逻辑学家和程序员关心的类型约定非常相似。

(2) 总体而言，逻辑学家更愿意不用带类型的变量。不过，没有类型时，集合论中就可能构造出某些悖论，导致无法回避的矛盾，使逻辑和数学推理崩溃。其中最著名的是罗素悖论：

\quad \hspace{-.4em}“设 $s$ 是所有\textit{不}以自身为元素的集合所组成的集合。

\quad $s$ 是否以自身为元素？”

\noindent
结果是，无论回答是还是否，都能立即证明你的回答错误。

罗素对这一悖论的解决办法，是为每个逻辑变量或数学变量关联一个\textit{类型}，规定它是个体、集合、集合的集合等等。然后规定，形如“$x$ 是 $y$ 的元素”的命题，只有在 $x$ 是个体类型的变量而 $y$ 是集合类型的变量时，或 $x$ 是集合类型而 $y$ 是集合的集合类型时，依此类推，才在语法上有意义。违反这一规则的命题都被视为无意义\textemdash{}根本谈不上真假，它只是一堆杂乱的字母。因此，涉及集合是否以自身为元素的任何命题，都可以直接排除。

罗素的类型论给数学基础带来了一些复杂性，这里无须讨论。对于我们的目的，它值得关注的特点是：类型被用来防止某些错误表达式进入逻辑和数学公式；而且，只需扫描文本就能检查是否违反类型约束，不必知道某个符号碰巧代表什么值。

(3) 在高级程序设计语言中，类型概念具有核心地位。同样，每个变量、常量和表达式都关联着唯一的类型。在 ALGOL 60 中，通过变量声明把类型与变量关联起来；在 FORTRAN 中，则根据变量名的首字母推断类型。在语言实现中，类型信息决定变量值的表示方式，以及必须为它分配多少计算机存储空间。类型信息还决定如何解释算术运算符，并使编译器能够把调用了不适当操作的程序作为无意义的程序拒绝接受。

由此可见，数学家、逻辑学家和程序员在使用类型概念时，有高度的共通性。类型概念的主要特点可以概括如下：

\begin{enumerate}[wide, nosep, label=(\arabic*)]
	\item 类型决定变量或表达式可以取值的类别。
	\item 每个值属于且只属于一个类型。
	\item 任何常量、变量或表达式所表示的值的类型，都可以根据其形式或上下文推断出来，无须知道它在运行时计算得到的值。
	\item 每个运算符要求操作数具有某种固定类型，并产生某种固定类型的结果（通常与操作数相同）。同一个符号用于多个不同类型时（例如 $+$ 既表示整数加法，也表示实数加法），可以把它视为一个有歧义的符号，代表几个不同的实际运算符。这种系统性的歧义总能在编译时消除。
	\item 一个类型的值的性质，以及定义在这些值上的基本操作的性质，由一组公理规定。
	\item 在高级语言中，类型信息既用于预防或检测程序中无意义的构造，也用于确定在计算机上表示和操纵数据的方法。
	\item 我们关心的类型，都是数学家已经熟悉的类型：笛卡尔积、带判别标记的联合、集合、函数、序列和递归结构。
\end{enumerate}

\subsection{数据类型的定义}

我们的数据结构化理论规定了若干定义类型的标准方法，以及如何在变量声明中使用这些类型，以规定变量在程序执行过程中可以取值的范围。大多数情况下，新类型用先前定义的\textit{组成类型}来定义。这种新类型的值是数据结构，可以由组成类型的\textit{分量}值构造出来，随后又可以从中提取分量值。这些分量值属于定义该结构化类型时使用的组成类型。如果只有一个组成类型，就称它为\textit{基}类型。

一个数据类型所含不同值的个数，称为它的\textit{基数}。许多类型的基数是有限的；用有限组成类型定义的结构化类型，其基数通常也是有限的，而且可以用简单公式计算。另一些数据类型的基数是无限的，例如整数；但它至多是可数无限的。原因在于，类型中的每个值都必须能通过有限次计算机操作构造出来，并且能在有限的存储空间中表示。任意实数、定义域为无限集的函数，以及其他具有不可数基数的类别，都不可能作为存储数据在计算机内部表示。不过，某些情况下可以用过程、函数或其他程序结构表示它们。

显然，一个结构的最终分量必然是非结构化的，而结构化类型的最终组成类型必然是非结构化类型。定义非结构化类型的一种方法，是简单地枚举其所有值，下一节将介绍这种方法。但某些情况下，最好把非结构化类型的性质视为由公理定义，并假定计算机硬件或高级程序设计语言的实现会把它们作为\textit{基本}类型提供。例如，ALGOL 60 的基本类型是 \textbf{integer}、\textbf{real} 和 \textbf{Boolean}；以下将假定这些类型已经可用。

\subsection{数据操纵}

数据在实践中最重要的方面，是如何操纵它，以及可用于操纵它的基本运算符有哪些。因此，我们为每个类型关联一组\textit{基本}运算符，希望它们既有助于程序设计，又至少存在一种效率合理的计算机实现。当然，基本运算符的选择多少带有任意性，所选的集合本来可以更大，也可以更小。指导原则是选出足够大的一组，使程序员需要的任何附加操作都能用这组基本操作定义，而且这样定义后也能高效实现。因此，如果某个运算符的高效实现方法在很大程度上取决于所选的数据表示方式，就把它看作基本运算符。

对于任意类型的数据，最重要、最一般的操作是赋值和相等测试。从概念上说，赋值把一个数据值从计算机存储器的一处完整地复制到另一处；相等测试则完整地扫描两个值（通常存储在不同位置），以检查它们是否相同。这些规则适用于基本数据类型，对结构化类型也没有理由偏离它们。如果结构化类型的值很大，这些操作可能花费相当长的时间；有时可以通过适当选择表示方式来缩短时间，也可以在把抽象程序转换为具体程序的过程中避免或消除这类操作。

另一类通用运算符是\textit{转换}函数，它们把一种类型的值映射为另一种类型的值。其中尤其重要的是\textit{构造函数}（constructor），它允许根据构成结构化类型的那些组成类型的值，定义该结构化类型的值。反方向的转换函数称为\textit{选择函数}（selector）；它们允许访问结构化类型的分量值。许多情况下，我们把已定义类型的名称用作标准构造函数或转换函数的名称，而该函数的值域就是这个类型。

某些数据类型适合看作有序的，并提供比较运算符来比较其值。但对于许多类型，排序没有任何有意义的解释，因此从抽象角度看，最好把它们视为无序的。这样有时会给以后具体设计中的表示方式和处理顺序策略的选择，带来更大的自由。

对于大型数据结构，高效操作的标准方法不是给它赋一个全新的值，而是\textit{选择性地更新}其中相对较小的一部分。通常的记法，是在赋值号左边写一个表达式（变量），用选择函数指出结构中要修改的位置。不过，我们也引入一些专门的赋值运算符，它们总以冒号开头，用于表示其他更一般的更新操作，例如向集合添加一个元素，或在序列末尾添加一个项。对于这两类选择性更新，都必须记住：从概念上或抽象角度看，即使只更新变量值的最小部分，整个变量值也已经改变。

\subsection{表示}

程序设计中的一个根本决定，是在多大程度上把计算结果存成数据以供以后使用，又在多大程度上按需计算。同样根本的决定，是如何在计算机中表示已存储的数据。许多简单、规模较小的情况下，都有一种显而易见的\textit{标准}表示方式，可以保证不会占用过多存储空间，也不会在执行基本操作时花费过多时间。但如果数据量（或处理量）很大，选择某种非标准表示往往有利，有时甚至是必需的。选择时要考虑所用存储介质（磁鼓、磁盘或磁带）的特性，以及各种操作的相对执行频率。表示细节的决定通常必须先于操纵数据的代码设计，并影响后者；而作决定时，数据的性质和所需的处理往往还不太清楚。因此，数据表示设计中很容易出现严重的判断错误，直到程序即将投入运行，甚至投入运行以后才暴露出来。这时，错误已经极难纠正。不过，在数据结构化中运用抽象，可以帮助我们推迟某些数据表示的决定，等到更加了解程序行为和数据特性时再作选择，从而减少这类错误，也使它们更容易纠正。

一个重要决定，是数据在存储时应当压缩到什么程度、以什么方式压缩，既节省空间，也节省输入$/$输出、复制操作和比较操作的时间。代价通常是其他所有操作所需的时间和代码量增加。比标准表示占用更少空间的表示，通常称为\textit{紧缩}表示；紧缩有不同程度，从松到紧不等。理论上值得关注的是\textit{最小}表示，它占用尽可能少的空间。在这种表示中，类型的值表示为从 0 到 $N - 1$ 的二进制整数，其中 $N$ 是类型的基数。对于基数无限的类型，如果每一种可能的位模式都表示该类型的一个值，那么这种表示就是最小表示。最小表示并不常用，因为处理它们的代价很高。

另一种节省空间的方法，是采用\textit{间接}表示。在数据的标准直接表示中，每个类型变量都分配足够的空间，以容纳该类型的任意值。间接表示中的变量，只需大到能容纳一个机器地址；在任一时刻，这个地址指向一个或多个机器存储单元，其中保存着当前值。当类型的基数无限时，这种技术是必需的，因为所需存储量会变化，而编写访问变量的代码时又无法知道这一存储量。如果实际所需存储量可变，并且在程序运行的大部分时间中都明显小于最大值，那么间接表示也可能有利。最后，如果认为许多不同变量往往会具有相同的值，也可以使用间接表示：只需保存该值的一份副本，各变量只保存指向它的指针；复制值的成本也很低，因为只需复制指针。但绝不能对这种共享副本进行选择性更新。

遗憾的是，间接表示往往需要动态存储分配和垃圾回收机制，因而带来额外的开销和复杂性；当数据必须在主存与辅助存储器之间复制时，还可能引发一些严重问题。

本章只介绍可能的数据表示中一小组有用的方法，熟练的程序员很容易再作补充。许多情况下，可以用一个更复杂但效率更高的数据类型定义，构造某种抽象数据类型的表示。例如，一个大集合可以表示为某个合适类型的项所组成的序列。后面各节会给出这样的例子。
\section[非结构化数据类型]{非结构化数据类型}
\label{sec:unstructured-data-types}

所有结构化数据，归根结底都必须由非结构化的分量构成，而这些分量属于基本类型或非结构化类型。其中某些非结构化类型（例如实数和整数）可以视为由程序设计语言或计算机硬件提供。虽然从理论上说，这些基本类型足以应付一切用途，但从实践看，有充分理由鼓励程序员定义自己的非结构化类型：既能明确说明变量可能的取值范围及每个值的含义，又能为以后设计高效表示创造条件。

特别是，许多计算机程序使用整数，并非表示数值数量，而是表示从少量备选项中选定的一项。这种情况下，程序注释通常会列出所有可能的值，并说明每个值的预期含义。可以把这样的量看作属于一个独立类型，它与整数类型截然不同，也与含义不同的其他类似标志集合截然不同。这种类型称为\textit{枚举类型}。我们提出一种标准记法，用来声明类型名，并为它的每个备选值关联一个名称：
\medskip

\tabto{3.7em}\textbf{type} suit $= ($club, diamond, heart, spade$)$;

\noindent
\hangindent=5.6\parindent\textbf{ordered} \tabto{3.7em}\textbf{type} rank $= ($two, three, four, five, six, seven, eight, nine, ten, Jack, Queen, King, Ace$)$;
\medskip

\tabto{3.7em}\textbf{type} primary color $= ($red, yellow, blue$)$;

\noindent
\hangindent=7.4\parindent\textbf{ordered} \tabto{3.7em}\textbf{type} day of week $= ($Monday, Tuesday, Wednesday, Thursday, Friday, Saturday, Sunday$)$;
\medskip

\tabto{3.7em}\textbf{type} day of month $= 1\,.\,.\,31$;

\noindent
\hangindent=6.2\parindent\textbf{ordered} \tabto{3.7em}\textbf{type} month $= ($Jan, Feb, March, April, May, June, July, Aug, Sept, Oct, Nov, Dec$)$;
\medskip

\tabto{3.7em}\textbf{type} year $= 1900\,.\,.\,1969$;

\tabto{3.7em}\textbf{type} Boolean $= ($false, true$)$;
\medskip

\noindent
\hangindent=5.6\parindent\textbf{ordered} \tabto{3.7em}\textbf{type} floor $= ($basement, ground, mezzanine, first, second$)$;

\tabto{3.7em}\textbf{type} coordinate $= 0\,.\,.\,1023$;
\medskip

前两个例子来自扑克牌。第一个声明指出，club、diamond、heart 和 spade 都是花色；换言之，任何 suit 类型的变量或表达式都只能表示这四个值之一，而标识符“club”“heart”“diamond”和“spade”就是该类型的\textit{常量}。类似地，\textit{rank} 类型的定义列出了十三个常量，表示该类型的十三个可能值。这种情况下，自然可以把类型视为有序的。接下来的例子声明了原色和一周各天的名称。考虑一个月中的日期时，把三十一个可能值全部写出很不方便。因此，我们引入如下约定：$a\,.\,.\,b$ 表示从 $a$ 到 $b$、包含两个端点的有限取值范围。它称为 $a$、$b$ 所属类型的\textit{子界}，这里是整数类型的子界。year 的声明也使用了这一约定。其他枚举类型的例子有：

\quad Boolean 类型，只有 false 和 true 两个值。

\quad month 类型，有十二个值，按所需顺序列出。

\hangindent=2.8\parindent\quad coordinate 类型，取值从 0 到 1023，可能用来表示阴极射线管显示器上的一个坐标。

以合适的方式定义类型后，程序员会用类型名指定变量的类型。为此，可以沿用数学家目前的惯例，把类型名写在变量\textit{后面}，两者以冒号分隔：

\quad trumps: suit; today:day of week;

\quad $pc$: primary color;

\noindent
如果同时声明几个同类型变量，可以采用简写形式，只列出变量名，不重复类型名：

\quad arrival, departure: day of month;

\quad $x$, $y$, $z$: coordinate.

\noindent
如果某个类型只用于少数变量，直接用类型定义替代类型名，会更加方便：

\quad answer: (yes, no, don't know);

枚举定义的类型，其基数显然等于定义列表的长度；数值子界的基数，则等于两个端点之差再加一。

\subsection{操纵}

要有效地操纵枚举类型和子界类型的值，需要以下操作：

\begin{enumerate}[wide, nosep, label=(\arabic*)]
	\item 相等测试，例如：

	\quad \textbf{if} arrival $=$ departure \textbf{then go to} transit desk;

	\quad \textbf{if} trumps $=$ spade \textbf{then} \dots

	\item 赋值，例如：

	\quad $pc \coloneq$ yellow;

	\quad trumps $\coloneq$ club;

	\item \label{en:manipulation-3}
	分情况选择，例如：

	\quad \textbf{case} $pc$ \textbf{of} $($red: \dots,

	\tabto{5.7em}yellow: \dots,

	\tabto{5.7em}blue: \dots$)$
	
	\noindent
	其中，$pc$ 是 primary color 类型的变量或表达式，省略号表示各个选择分支。分情况选择可以是一条语句，此时各分支必须是语句；也可以是一个表达式，此时各分支必须是同一类型的表达式。
	
	分情况选择的效果，是选出唯一一条语句（或一个表达式）来执行（或求值）：该分支前面的常量等于 case 表达式的当前值。有时，把几个常量都放在同一分支前面会比较方便；也可以指定一个子界，表示其中的值都会选中相应分支。当然，每个值都必须且只能出现一次：
	
	\quad \textbf{case} digit \textbf{of} ($0\,.\,.\,2:\dots$,
	
	\tabto{6.6em}$3:7:\dots$,
	
	\tabto{6.6em}$4\,.\,.\,6:8:9:\dots$).
	
	\noindent
	最后这个例子中，用基本字 \textbf{else} 替代最后一个分支的标号会比较方便，它涵盖前面各分支没有明确列出的所有剩余情况。
	
	选择分支为语句时，我们有时用花括号代替圆括号，把它们括起来。
	
	\item 对于声明为有序的类型，可以测试各个值之间的次序关系：
	
	\quad \textbf{if} May $\leqslant$ this month \& this month $\leqslant$ September \textbf{then}
	
	\qquad adopt summer timetables. 
	
	\noindent
	其他情况下，值的次序完全无关，对程序员没有意义。
	
	\item 对于有序类型，引入后继函数和前驱函数（succ 和 pred）很有用，它们把类型中的每个值映射为紧接着的较大值或较小值，前提是这样的值存在。此外，如果 $T$ 是任意有序类型，记法 $T$.min 表示其最小值，$T$.max 表示其最大值，前提是它们存在。这样有助于在表述程序、定理和公理时，不依赖常量的具体名称。
	\item 在计算机程序中，我们常常希望让变量依次遍历一个类型的所有值。这可以用下面这种
	
	for 语句或循环表示：
	
	\quad \textbf{for} $a$: alpha \textbf{do} \dots;
	
	\quad \textbf{for} $i$: $1\,.\,.\,99$ \textbf{do} \dots;
	
	\noindent
	在这种构造中，计数变量（$a$ 或 $i$）属于所指定的类型，并且在该构造内部局部声明；也就是说，它的值在循环之前或之后都不存在，循环外也不能访问它的名称。此外，不允许在循环体内改变计数变量的值，否则就违背了用 \textbf{for} 构造声明该变量的整个意图。
	
	对于有序类型，自然假定计数变量按定义的顺序遍历类型的各个值：$T.$min、succ($T.$min)、\dots、$T.$max。但对于无序类型，则假定遍历顺序在当前抽象层次上无关紧要，具体顺序留待开发具体程序的后续阶段确定。
	
	\item 对于子界类型，尤其是整数子界，有时需要进行原来较大类型中定义的操作。原则上，这很简单：先把子界值转换为较大类型中的对应值，然后进行操作，最后在必要时再转换回来。这需要类型转换函数；为此，使用目标类型的名称很方便，例如：
	
	\quad $x$distance $\coloneq$ integer$(x)\ -$ integer$(y)$;
	
	\quad $z \coloneq$ coordinate$($integer$(z) + x$distance$)$;
	
	\noindent
	其中 $x$distance 是整数变量。当然，这种记法过于繁琐；如果能根据上下文确定应该在哪里补回转换，就一定希望约定省略这些转换：
	
	$x$distance $\coloneq x - y$;
	
	$z \coloneq z + x$distance.
\end{enumerate}

\noindent
\textit{练习}

\noindent
给定 $m$:month 和 $y$:year，写一个分情况选择表达式，给出月份 $m$ 的天数。

\subsection{表示}

枚举类型 $T$ 的标准表示，是按声明的顺序，把各个值映射为从 0 到 $n - 1$ 的计算机整数，其中 $n$ 是类型的基数。因此，这种情况下标准表示也是最小表示。子界的标准表示，是让每个值沿用它在原类型中的表示；这样，在类型之间转换就不需要实际操作。当然，从基类型转换到子界类型时，应当检查该值是否在规定范围之内。

子界值的最小表示，可以从标准形式中减去子界最小值的整数表示而得到。这种情况下，转换为子界值时，除了检查还需要做减法；反向转换则需要做加法。

除这些转换外，无论采用哪一种表示，枚举类型和子界类型都可以以相同方式处理。次序测试可以用计算机通常的整数指令完成；succ 和 pred 则分别加一或减一，随后测试结果是否仍在范围内。

分情况选择可以用选择跳转最有效地实现。例如，上述第一个例子（3.1.\ref{en:manipulation-3}）在 ALGOL 60 中可以编码为：

\quad \textbf{begin switch} $ss \coloneq$ red, yellow, blue;

\quad \quad \textbf{go to} $ss[pc + 1]$;

\quad \quad red:\textbf{begin}\dots; \textbf{go to} end \textbf{end};

\quad \quad yellow:\textbf{begin}\dots; \textbf{go to} end \textbf{end};

\quad \quad blue:\textbf{begin}\dots; \textbf{go to} end \textbf{end};

\quad end:\textbf{end}. 

它可以高效地表示为机器码，使用带变址的跳转指令，以及列出各选择分支对应代码起始地址的选择表。

for 语句的实现，显然对应于步长为一的 ALGOL 60 for 语句。上面提出的约定，把计数变量视为循环的局部常量，既有助于文档表达清楚，也有助于利用寄存器、专门的计数和测试指令等，提高计算机执行效率。

\subsection{例子}
\label{sec:example-3-3}

计算机某一外部设备的字符集由枚举定义：

\quad \textbf{type} character $= (\dots)$;

\noindent
这个集合包含以下子界：

\quad \textbf{type} digit $=$ nought\,.\,.\,nine;

\quad \textbf{type} alphabet $= A\,.\,.\,Z$;

\noindent
还包含独立符号 point、equals、subten、colon、newline、space，以及若干其他单字符运算符和标点符号。

有一个变量：

\quad buffer:character 

\noindent
它保存从该外部设备最近输入的字符。通过过程“read next character”，可以从输入带向 buffer 输入一个新值。

在 ALGOL 60 的某一种表示中，基本字不用下划线标出，因此看起来像标识符。所以，如果其前面或后面紧邻一个标识符或数，就必须用一个或多个空格或换行符把两者分开。

ALGOL 翻译程序的第一遍处理，需要读入单个字符，把它们组合成语言中有意义的符号。因此，一个标识符、一个基本符号、一个数，以及赋值号“$\coloneq$”，各自都算一个符号；其他所有标点符号也是如此。空格和换行在完成分隔符号的作用后，就必须忽略。我们假定，每个有意义的符号都由一个专门设计的例程扫描，而每个例程都会在 buffer 中留下输入中第一个\textit{不}属于该符号的字符。

作为程序符号分析的一个例子，输入文本

\quad $l$:beta1 $\coloneq \text{beta} \times 12$;

\noindent
应当分析成以下符号：

\quad $l$

\quad :

\quad beta1

\quad $\coloneq$

\quad beta

\quad $\times$

\quad 12

\quad ;

程序的总体结构，是根据符号的首字符进行分情况选择，从而确定符号属于哪一类。

read first character;

\textbf{repeat case} buffer \textbf{of}

\quad $($alphabet:scan identifier,

\quad digit:point:subten:scan number,

\quad space:newline:read next character,

\quad colon:\textbf{begin} read next character;

\quad\quad \textbf{if} buffer $=$ equals \textbf{then}

\quad\quad\quad \textbf{begin} deal with ``becomes''; read next character \textbf{end}

\quad\quad \textbf{else} deal with single colon character

\quad\quad \textbf{end}

\quad \textbf{else begin} deal with single character;

\quad\quad read next character

\quad \textbf{end}

\quad $)$

\textbf{until} end of tape
\section[笛卡尔积]{笛卡尔积}

已定义的枚举类型和子界，与基本数据类型一样，原则上都是非结构化的。当然，这些类型的任何具体表示都有结构，例如一组连续的二进制位；但从抽象角度看，这种结构实质上无关紧要。我们不提供访问各个位或用各个位构造值的运算符。实际上，要成功运用抽象，就必须忽略这种可能性；唯有如此，才能推迟有关细节的决定，依据计算机的特性以及数据的操纵方式来决定数据表示。

现在转而讨论这样的数据类型：其结构对程序员有意义，至少在程序开发的某个阶段有意义。我们方法的基础是，和枚举类型一样，程序员应当能够通过声明引入新数据类型；但对于结构化数据，新类型的定义会引用其他基本类型或先前定义的类型，也就是结构分量的类型。因此，声明新类型有些类似于在 ALGOL、FORTRAN 这类语言中声明新函数。函数声明用现有的或先前声明的函数和操作来定义新函数。已声明的函数可以在程序语句或其他函数声明中多次调用；同样，新类型也可以在程序的其他声明中多次“调用”。这些声明可以是变量声明，规定变量在新类型中取值；也可以是另一个新类型的声明。

首先讨论初等数据结构，即笛卡尔积和联合。对于数学家和逻辑学家，这些初等结构几乎和自然数一样简单、熟悉。此外，从计算机程序员的角度看，只要组成类型也都是初等类型，初等数据结构就具有非常有利的性质。

\begin{enumerate}[wide, nosep, label=(\arabic*)]
	\item 首先，每个数据项占用固定、有限而且通常不多的磁芯存储空间；所需空间只随定义的大小线性增长。

	\item 保存每个值所需的存储空间，既可以在主存中永久分配，也可以在运行时栈上分配，两种方式都很高效。不需要更复杂的动态存储分配系统。

	\item 对数据项最有用的那些操纵，可以在当今计算机上用简单、紧凑的机器码指令序列高效完成。

	\item 这些结构的表示不需要指针（引用、地址），因而把这类数据在主存和辅助存储器之间传送时不会有问题。

	\item 对于任何给定结构，选择合适表示通常不会给程序员造成困难。
\end{enumerate}

我们首先讨论的数据结构化方法是笛卡尔积。一个熟悉的例子是复数空间：每个复数都由一对浮点数构成，其中一个看作实部，另一个看作虚部。复数类型的声明可以写成：

\quad \textbf{type} complex $= ($real part:real; imagpart:real$)$;
\nopagebreak

\noindent
或更简洁地写成：
\nopagebreak

\quad \textbf{type} complex $= ($realpart, imagpart:real$)$.

定义引入 realpart 和 imagpart 这两个名称，提供选择复数分量的手段。例如，若 $n$ 属于上述 complex 类型，$n$. realpart 就表示其实部，$n$. imagpart 表示其虚部。

表示笛卡尔积类型中某个值的常量，可以用表示各分量值的一列常量来定义。如前所述，类型名用作转换函数，以指出所构造结构的类型；它接受一列参数，而不是单个参数。因此，复数 $13 + i$ 可以写成：

\quad complex$(13, + 1)$.

\noindent
笛卡尔积的另一个例子，是声明一种类型，其值表示扑克牌。每张牌都可以先给出花色（例如 heart），再给出点数（例如 Jack）来指定。要唯一指定一张牌，这两项信息缺一不可。因此，cardface 类型可以定义为 suit 和 rank 类型的笛卡尔积：

\quad \textbf{type} cardface $= (s$:suit; $r$:rank$)$.

该类型的典型常量有：

\quad cardface$($club, two$)$, cardface$($heart, Jack$)$.

另一个简单的笛卡尔积例子是日期，这次它有三个
分量。通常，日期可以用三个
值来指定：第一个从 day of
month 类型的可能值中选出，例如第七天；第二个从
month 类型的可能值中选出，例如 March；第三个从
year 类型的值中选出，例如 1908。这个日期可以写成：

\quad date$(7$, March, $1908)$.

\noindent
它属于这样声明的类型：

\quad \textbf{type} date $= ($day:day of month; $m$:month; $y$:year$)$;

笛卡尔积类型的定义特征是，它包含各分量类型的值的一切可能组合，即使其中一些组合在实践中绝不会出现。所以，date $(31$, Feb, $1931)$ 是 date 类型的一个正常值，尽管现实世界中不存在这样的日期。但是，date $(28$, Feb, $1899)$ \textit{不是} date 类型的值，因为按上面的定义，1899 不是 year 类型的值。因此，date 类型的定义与现实情境并不完全一致，但对于大多数用途，两者已足够接近；程序员有责任确保，操纵这类变量时，绝不会使它们取到自己认为无意义的值。

这个例子表明，我们所提供的、用其他类型定义新类型的方法，比数学上用其他集合定义新集合的一般方法更简单，表达能力也更弱；因为数学上当然可以定义一个集合，排除所有不合要求的日期。事实上，在声明类型或变量时，严格说明每个有意义的值都将具有的性质，是编写文档的良好做法。
 
最后一个例子说明，如何把二维栅格上的点位置集合，声明为一维坐标的笛卡尔积：

\quad \textbf{type} raster $= (x$, $y$:coordinate$)$

这是用笛卡尔坐标从一维构造二维空间的标准方法；二维空间中的每一点，都可以用简单一维数的有序对来标识。这就解释了为什么用“笛卡尔积”这个词来称呼上述类型定义方法。如果 $r$ 是 raster 类型的变量，$r.x$ 和 $r.y$ 通常分别称为 $r$ 在 $x$ 轴和 $y$ 轴上的投影；不过，我们会把函数 $x$ 和 $y$ 称为选择函数，而不是投影。

笛卡尔积类型的基数，等于各组成类型的基数之积。如果把笛卡尔积想象成一个矩形或长方体，各边长度等于相应坐标轴类型的基数，就很容易理解这一点。因此，纸牌类型的基数是十三乘四，即五十二；正如预期，这就是一副标准扑克牌的张数。日期的个数是 $26\,040$，略多于相应时间区间内实际的天数，因为如前所述，其中包含少量无效日期。

\subsection{操纵}

除所有类型共有的赋值和相等测试外，积类型的主要操作，就是用分量值构造一个值，以及选择其中的分量。构造笛卡尔积类型的值时，原则上必须写出类型名，把它用作转换函数。但在不会引起混淆时，沿用传统数学做法，把转换函数省略，通常更方便。如果类型根本没有明确名称，就必须这样做。例如，可以写 $($heart, Jack$)$，而不写 cardface $($heart, Jack$)$。

选择分量时，我们使用点记法，例如 $n$. imagpart。这比通常的函数记法 imagpart$(n)$ 更方便，因为它避免了不必要的深层括号嵌套。另一个极重要的操作，是对变量的分量进行选择性更新。这可以把分量名写在赋值号左边来表示：

\quad $u$.imagpart $\coloneq$ 0;

\quad $r.x \coloneq a\times r.x + b\times r.y$.

如果把笛卡尔积声明为有序的，那么所有组成类型都必须有序；自然的做法是按字典序定义其顺序，排在前面的分量优先。因此，如果 suit 和 rank 都有序，就可以把 cardface 类型声明为有序，采用传统顺序：所有梅花在所有方块之前，随后是所有红桃和所有黑桃；而在同一花色内部，按点数排列。

检查或处理一个结构化值时，经常需要在一小段代码中多次引用它的分量。这时，使用 \textbf{with} 构造会很方便：

\quad \textbf{with} $sv$ \textbf{do} $S$;

\noindent
其中，$sv$ 指出该结构化变量（或表达式），$S$ 是规定如何处理它的程序语句。在语句 $S$ 中，只需用各分量的选择函数名称 $s_1, \dots, s_n,$ 来引用 $sv$ 的分量，不必使用通常的写法：$sv.s_1, sv.s_2, \dots sv.s_n$。使用这种构造的理由是：

\begin{enumerate}[wide, nosep, label=(\arabic*)]
	\item 使这一段程序的目的更加明确。
	\item 简化它的表述。
	\item 指出实现时可能提高效率的机会。
\end{enumerate}

\noindent
\textit{例子：}给定 today:date，测试它是否是有效日期。

\quad \textbf{with} today \textbf{do case} $m$ \textbf{of}

\quad \quad $\{$Sept:April:June:Nov:

\quad \quad \quad \textbf{if} day $>$ 30 \textbf{then go} to invalid,

\quad \quad \quad Feb:\textbf{if} day $> (\text{\textbf{if }} (y \div 4) \times 4 = y \text{ \textbf{then} } 29 \text{ \textbf{else} } 28)$

\quad \quad \quad \quad then go to invalid,

\quad \quad \quad else do nothing$\}$.

\noindent
\textit{练习}
\nopagebreak

\noindent
编写函数，表示复数的四种标准算术运算。

\subsection{表示}
\label{sec:representation-4-2}

表示笛卡尔积类型的值，标准方法就是把分量值并排放在一段连续的存储区域中，通常遵循定义给出的顺序。不过，这种并排放置涉及的紧缩和填充程度，可以有很大差别。在标准的非紧缩表示中，每个分量值都占用整数个字；字是计算机中便于寻址且能高效访问的最小存储单位。

如果一个值所需的空间不足一个字，就可以把多个分量紧缩在一个字中。在紧密的紧缩表示中，各分量的位模式直接拼接在一起。在较松的紧缩表示中，分量可能放在字的某些分区内；这些分区是“自然的”，因为计算机有专门的机器码指令来选择或更新一个字的特定部分\textemdash{}例如字符边界，或字中的指令字段。

可以重新排列分量顺序，以便把它们安置在这样的边界内；但如果类型有序，这种重排通常并不可取。

如果紧缩表示跨越多个字，单个分量值就可能跨越字边界。在许多机器上，选择或更新这种分量会比通常情况耗费多得多的时间。因此，常见做法是在字末留下少量未使用空间（填充），避免这种跨界。

要构造结构化值的最小表示，必须对所有分量都采用最小表示。然后，把每个分量乘以它后面所有分量类型的基数之积，再把这些结果相加，就得到笛卡尔积类型的最小表示。例如，1908 年 3 月 7 日的表示是 $6\times 12\times 70 + 2\times 70 + 8 = 5188$。

各种表示之间的选择，取决于处理这些值的更广泛背景。如果选择和选择性更新频繁，采用非紧缩表示就比较划算，因为这些操作可以直接使用按字寻址硬件的常规选择机制。但如果比较频繁的是整体值的复制和比较，那么采用紧缩表示就更有利，因为完成这些操作需要的指令和存储访问次数更少。还应考虑一种特殊的复制，即计算机主存与辅助存储器之间的传送。如果这种传送频繁，恰当紧缩以减少传送数据量，可以显著提高效率。

\begin{figure}[th]
	\centering
	\begin{subfigure}[t]{0.3\textwidth}
		\centering
		\includegraphics[height=4em]{chap2/fig1a}
		\caption{}
	\end{subfigure}
	~
	\begin{subfigure}[t]{0.3\textwidth}
		\centering
		\includegraphics[height=4em]{chap2/fig1b}
		\caption{}
	\end{subfigure}
	~
	\begin{subfigure}[t]{0.3\textwidth}
		\centering
		\includegraphics[height=4em]{chap2/fig1c}
		\caption{}
	\end{subfigure}
	\caption{日期 (7, March, 1908) 的表示}
\end{figure}

采用紧缩表示的第二种情况，是数据存储空间不足，无论是在主存还是外部辅助存储器中。不过，每当检查或更新数据时，都需要解开紧缩、再重新紧缩，这会使代码膨胀；必须注意，节省的数据存储空间不能被增加的代码空间抵消。

最小表示不常用于数据存储，因为它额外节省的空间很少（每个分量总是不到一位）；与通常的移位和掩码操作相比，每次访问分量都要乘除，增加的时间通常远远超过这点收益。不过，如果没有其他办法能把值放在合适的字边界之内，这种技术就可能有用，也可以与较常规的紧缩方法结合使用。此外，如果这个值只用于或主要用于给多维数组作索引，就应优先采用最小表示，因为这样能显著节省数组表示所需的空间（见第 \ref{sec:representation-6.2} 节）。

把 \textbf{with} 构造表示为机器码时，有时可以先计算所引用结构的地址，并把它保存在寄存器中；这样，访问分量的代码可能更短、更快。如果分量已经紧缩，可以在开始处理之前把它们解开，放入独立的字，以方便引用和更新；如果更新过这些分量，就必须在处理结束时重新紧缩，再存回结构中。在某些机器上，一次紧缩或解开整个结构，比逐个分量进行这些操作更经济。

\noindent
\textit{练习}
\nopagebreak

\noindent
给定变量

\quad today:date;

\noindent
编写程序，把紧接其后的日期赋给变量 tomorrow: date。采用紧密的紧缩表示，把这个程序翻译为你选择的机器码。分别使用非紧缩表示和最小表示重写程序。比较各自代码的长度，以及执行所需的时间。
\section[带判别标记的联合]{带判别标记的联合}

定义对象集合时，把一个集合定义为两个已知集合的并集，往往很有用。例如，在一副标准扑克牌中加入大小王后，扩展后的集合可以描述为标准牌集合与“百搭牌”joker 1、joker 2 所组成集合的并集。以这个集合的元素为取值的类型，可以声明为两种备选类型的联合：纸牌类型，以及有两个不同值的枚举类型：

\quad \textbf{type} pokercard $= ($normal:$(s$:suit; $r$:rank$)$,

\tabto{10.2em}wild:$($joker 1, joker 2$))$.

\noindent
pokercard 类型的每个值，要么对应一个有序对，其分量表示花色和点数；要么对应枚举类型中的一张王。

指定带判别标记的联合类型的常量时，必须指出它所表示的值来自哪一种备选类型。为此，要明确写出备选项的名称，例如：

\quad pokercard$($normal$($heart, Jack$))$

\noindent
表示来自第一种备选项的值，而

\quad pokercard$($wild$($joker 2$))$

\noindent
表示来自第二种备选项的值。一般而言，如果能从上下文推断类型，省略类型名会比较方便。

带判别标记的联合的第二个例子，可以来自维护一国所有汽车登记册的工作。汽车可以分为本国居民拥有的本地汽车，以及非居民临时带入境内的访客汽车。这两种情况下所需的信息很不相同。两种汽车的车牌号码和厂牌都被认为是有关信息。但对于本地汽车，还需要车主姓名，以及汽车首次登记在该车主名下的日期。对于访客汽车，这些信息无关紧要，只需要来源国名称的标准三字母缩写。因此，这两种备选汽车类型可以定义为：

\quad \textbf{type} local car = (make:manufacturer; regnumber:carnumber;

\tabto{9.4em} owner:person; first registration:date);

\quad \textbf{type} visitor car = (make:manufacturer; regnumber:carnumber;

\tabto{10.1em} origin:country);

现在就可以定义一个涵盖两类汽车的类型：

\quad \textbf{type} car $=$ (local:local car,

\tabto{7.4em} foreign:visitor car).

\noindent
不过，分别定义本地和外国汽车的结构并不方便；我们希望利用两者有若干相同分量这一事实。可以把共有分量移到两个备选项之前：

\quad \textbf{type} car $= ($make:manufacturer;

\tabto{7.5em}regnumber:carnumber;

\tabto{7.5em}$($local:$($owner:person;

\tabto{7.8em}first registration:date$)$,

\tabto{7.8em}foreign:$($origin:country$))$

\tabto{7.1em}$)$.

\noindent
每辆汽车都有 make 和 regnumber，但只有本地汽车才有 owner 和 first registration 日期；只有外国汽车才有 origin。

第三个例子是几何图形的定义。在某个应用中，图形可以分为矩形、三角形和圆：

\quad \textbf{type} figure $= ($position:point; rect:$R$, tri:$T$, circ:$C)$.

\noindent
各类图形的指定方式不同。对于矩形，要给出一条边的倾角，以及两条边的长度：

\quad \textbf{type} $R = ($inclination:angle; side1, side2:real$)$.

\noindent
三角形用一条边的倾角和长度，以及这条边与另外两条边形成的角来指定：

\quad \textbf{type} $T = ($inclination:angle; side:real; angle1, angle2:angle$)$.

\noindent
对于圆，只需用一个实数指定直径：

\quad \textbf{type} C $= ($diameter:real$)$.

把类型定义为若干其他类型的联合时，必须认识到：它的值应当被视为与任何组成类型的值完全不同。否则，就会立即违反“每个值只属于一个类型”的规则。因此，类型的联合必须与通常的集合并集概念明确区分。此外，对于联合类型的每个元素，都能确定它来自哪一种组成类型，即使同一种类型重复出现了多次。例如，玩纸牌接龙所用的两副牌，可以定义为两个牌组的联合：

\quad \textbf{type} patience card $=$ (red:cardface, blue:cardface).

\noindent
patience card 类型的每个值都明确标记着，它来自红色牌组还是蓝色牌组，即使现实中牌背颜色可能相同。这就解释了为什么把这种类型定义形式称为“带判别标记的联合”。由此可知，带判别标记的联合的基数，总等于各组成类型的基数之和。

\subsection{操纵}

带判别标记的联合的任何值，都附带一个\textit{判别标记}（tag）字段，指出它来自哪一种具体组成类型。赋值时，这个字段也要复制；相等测试时，要使两个值相等，其判别标记字段必须相同。

构造带判别标记的联合类型的值时，必须写出该值来源的备选类型名称：
\smallskip

\quad patience card$($red$($spade, Jack$))$.
\smallskip

\noindent
这样就会自动把值“red”赋给结果的判别标记字段。

某辆特定汽车可以表示为：
\smallskip

\quad car$($Ford, ``RUR157D'', local$($me, date$($1, Sept, 1968$)))$.
\smallskip

要访问并操作编码为带判别标记的联合的信息，需要把它转换回原来的类型。可以约定把该类型的标号当作选择函数使用，例如：

card1.wild \tabto*{8em} 的类型是 $($joker 1, joker 2$)$

car1.foreign \tabto*{8em} 的类型是 $($origin:country$)$

fig1.tri \tabto*{8em} 的类型是 $T$

\noindent
如果组成类型是笛卡尔积，就可以对所得值合法地使用它的选择函数；约定点运算符向左结合。

\quad card1.normal.$r$

\quad car1.local.owner

\quad fig1.circ.diameter

如果程序员试图把带判别标记的联合的值转换回一个\textit{并非}其来源的类型，就犯了严重的程序设计
错误，可能产生无意义的结果。这种错误只能通过运行时检查来发现：每当显式或隐式调用这种转换时，都检查判别标记字段。检查既费时，一旦失败又很不方便。因此，我们希望找到一种记法，保证这种错误绝不会在运行中的程序里发生；而且，只检查文本就能作出保证，无须知道程序运行时正在处理的值。如果有自动编译器，也可以由它提供这样的保证。

我们提出的记法，是笛卡尔积的 \textbf{with} 构造与分情况选择的 case 构造的结合。设联合类型的值 $sv$，需要根据它来自哪种备选类型，以几种不同方式之一进行处理。那么可以写：

\quad \textbf{with} $sv$ \textbf{do} $\{a1:S_1$

\tabto*{7.6em}$a2:S_2$,
\smallskip

\tabto*{8.85em}\vdots

\tabto*{7.6em}$an:S_n\}$;

\noindent
其中，每当 $sv$ 的判别标记字段值为 $a_i$ 时，就选择执行语句 $S_i$。在 $S_i$ 内，只要 $sv$ 的值保持不变，就可以安全地假定它来自对应的备选类型。因此，可以直接使用为该备选类型定义的分量选择函数来引用 $sv$ 的分量，就像前面为笛卡尔积介绍的简单 \textbf{with} 语句一样。

如果希望把联合类型看作有序的，最自然的顺序是：列表中较早备选项对应的所有值，都排在较晚备选项的任何值之前。

\noindent
\textit{练习}
\nopagebreak

\noindent
编写函数，计算按上述定义给出的图形的面积。

\subsection{表示}

表示带判别标记的联合的值时，首先需要把判别标记表示为从零到 $n - 1$ 的一个整数，其中 $n$ 是备选类型的个数。判别标记后面紧接着原类型的值的表示。与笛卡尔积一样，表示的紧缩程度可以选择。

在非紧缩表示中，判别标记占一个完整的字；联合类型的每个值所占空间，比最大备选类型的值多一个字。在紧缩表示中，这一额外开销可以缩减为几位。在最小表示中，把每个值在原类型中的最小表示，加上联合中所有排在它之前的类型的基数之和，就得到该值的表示。因此，来自第一种类型的值，例如 (diamond, four)，其表示值与原类型中的表示完全相同，即 16。但 joker 1 在原枚举类型中的表示值为零，需要再加上纸牌类型的基数。

在非紧缩、紧缩和紧密的紧缩表示之间选择时，所依据的考虑与笛卡尔积相同。不过，最小表示的运行速度代价要小得多，因为恢复原值只需减法，不需除法。

一般而言，不同备选类型的值占用的存储空间不同，因此必须为较短的值“补齐”填充，使各值长度相同，从而遵守初等数据类型占用固定存储量这一方便的规则。后面的章节将会说明，当值是某个更大结构的分量时，这种填充常常可以省去。

\begin{figure}[ht!]
	\centering
	\begin{subfigure}[b]{0.4\textwidth}
		\centering
		\includegraphics[height=6em]{chap2/fig2a}
	\end{subfigure}
	~
	\begin{subfigure}[b]{0.4\textwidth}
		\centering
		\includegraphics[height=6em]{chap2/fig2b}
	\end{subfigure}	
	\caption{汽车的表示}
\end{figure}

在当今程序设计实践中，联合的表示省略判别标记字段，是相当常见的做法。要正确操作这种表示，程序员必须根据其他因素“知道”该值应如何解释，因为无法从值本身查明。如果他的判断有误，无论运行时还是编译时检查，都无法发现。由于这类错误的效果取决于位模式表示的细节，所得结果无法用程序员正在使用的抽象概念来预测。因此，一般看来，应当把使用判别标记字段、采用可在编译时检查的分情况选择作为标准做法，只有特殊情况下才绕过它们。

\subsection{例子}

回到第 \ref{sec:example-3-3} 节的例子，即把语言文本分析为有意义的符号。我们希望严格地抽象定义这些符号究竟是什么。

\quad \textbf{type} symbol $=$

\tabto*{5.6em}$($realconst:real,

\tabto*{6em}integerconst:integer,

\tabto*{6em}identifier:ident,

\tabto*{6em}basic:delimiter$)$;

\noindent
其中，暂时不定义 ident 类型，并假定各个界符通过枚举定义。
\section[数组]{数组}
\label{sec:array}

对于许多程序员，数组是最熟悉的数据结构；在某些程序设计语言中，它还是唯一明确提供的结构。从抽象角度看，数组可以视为一个映射：从某个类型的定义域（下标范围），映射到一个可能属于不同类型的值域（数组的类型，或更准确地说，数组元素的类型）。

数学家通常用箭头指定映射的类型：

\quad $M:D\rightarrow R$;

\noindent
其中 $D$ 是定义域类型，$R$ 是值域类型。另一种程序员更熟悉的记法是：

\quad M:\textbf{array} $D$ \textbf{of} $R$.

\noindent
这种记法更能体现数据的表示方式，而数学记法强调的是结构的抽象性质，与具体表示无关。

把映射类型的某个具体值 $M$ 应用于定义域类型的值 $x$ 时，它指定值域类型中的唯一一个元素，称为 $M$ 在 $x$ 处的值，可以用圆括号或方括号写成：

\quad $M(x)$ 或 $M[x]$。

\noindent
映射也叫函数。这里采用“映射”一词，是为了把这种数据结构与一段程序区分开：后者确实根据定义域中的参数，计算出值域中的一个值。区别的本质在于，映射 $M$ 并非通过给出计算方法来指定，而是通过明确列出定义域中每个可能值 $x$ 所对应的 $M(x)$ 来指定。因此，数组只能用于在有限个点上定义的函数，而计算函数的定义域可以是无限的。

有限映射的一个例子是 monthtable，它为一年中的每个月指定该月的天数：

\quad \textbf{type} monthtable $=$ \textbf{array} month \textbf{of} $28\,.\,.\,31$.

\noindent
定义域是 month 类型，值域类型由 28 到 31（包含端点）的整数构成。要指定这个类型的一个典型值，只需让 $x$ 遍历定义域，并列出各个 $M(x)$ 的值。因此，如果 $M$:monthtable 被指定为：

\quad monthtable$($Jan:31, Feb:28, March:31, April:30,

\quad May:31, June:30, July:31, Aug:31,

\quad Sept:30, Oct:31, Nov:30, Dec:31$)$

\noindent
那么 $M[\text{Jan}] = 31,\, M[\text{Feb}] = 28$，依此类推。

数组提供了一种表示一副牌中某种具体排列的方法，因为每种排列都可以视为一个映射，为牌组中五十二个可能位置各自指出占据该位置的纸牌。因此，每种可能的排列都可以视为以下映射类型的一个值：

\quad \textbf{type} cardpack = \textbf{array} $1\,.\,.\,52$ \textbf{of} cardface.

\noindent
当然，并非该类型的所有值都表示实际的牌组，因为没有任何条件阻止某个值把两个不同位置映射到同一张牌，而现实中这是不可能的。

元素属于笛卡尔积类型的数组，有时称为\textit{表}。

数组的第三个例子，表示普通穿孔卡上字符穿孔的一切可能配置。它可以看作映射 $M$，把每个列号映射为一个字符，即该列穿孔所表示的字符。

\quad \textbf{type} punchcard = \textbf{array} $1\,.\,.\,80$ \textbf{of} character.

\noindent
卡片上任何可能的穿孔文本，都可以视为 punchcard 类型的一个值。

第四个例子用数组表示阴极射线管显示设备上一页内容的某个可能值。假定它有 40 行，每行有 27 个字符位置。把映射的定义域指定为所有可能行号与每行所有可能字符位置的笛卡尔积，就可以得到二维效果。写法如下：

\quad \textbf{type} spot $= ($row:$1\,.\,.\,40$; column:$1\,.\,.\,27)$;

\quad \textbf{type} display page $=$ \textbf{array} spot \textbf{of} character.

处理多维数组的另一种方法，是把它看作由行组成的数组，而每一行又是字符数组：

\quad \textbf{type} display page $=$ \textbf{array} $1\,.\,.\,40$ \textbf{of} row;

\quad \textbf{type} row $=$ \textbf{array} $1\,.\,.\,27$ \textbf{of} character.

\noindent
如果要逐行处理而不逐列处理，这种抽象结构更合适。

数组类型的基数，是值域类型的基数的若干次幂，指数为定义域类型的基数，即：

\quad cardinality$(D \rightarrow R) =$ cardinality$(R)^{\text{cardinality}(D)}$

\noindent
可以通过考察完整指定一个数组类型的值需要作多少次选择，来证明这一点。对于定义域的每个值，都要从值域类型的 cardinality$(R)$ 个可能值中选出一个。每个数组元素都要独立地作一次这样的选择，也就是总共选择 cardinality$(D)$ 次。

\subsection{操纵}

把定义域的所有值都映射到值域中同一个值的映射，称为常值映射。数组的一种自然构造函数，接受任意值域值作为参数，产生一个常值数组，其中所有元素都等于给定的值域值。直接用类型名来表示这个构造函数会很方便，例如：

\quad $M = \text{monthtable}(31)$

\noindent
是一个数组，对于所有月份 m，都有 $M[m] = 31$。

\quad card pack$($cardface$($heart, King$))$

\noindent
显然是一副魔术师用的牌。

数组的基本构造操作，是为数组中某个特定元素定义一个新值。如果 $x$ 是数组类型 $T$ 的值，$d$ 是其定义域类型的值，$r$ 是其值域类型的值，那么写：

\quad $T(x, d:r)$

\noindent
表示 $T$ 类型的一个值，它除把 $d$ 映射到 $r$ 之外，在所有方面都与 $x$ 相同。如果能从上下文推断出 $T$，就可以省略它。类似地，常值数组 $T(x)$ 可以写为 all $(x)$。

数组的基本选择运算符是下标运算。它实质上是一个二元操作，作用于数组和一个定义域类型的值，产生对应的值域类型的值。

改变数组值最常见、最高效的方法，是选择性地更新某个分量。采用通常的记法，把带下标的数组变量放在赋值号左边，就能实现这一点：

\quad $a[d]\coloneq r.$

\noindent
它与下式意义相同：

\quad $a\coloneq T(a, d:r).$

\noindent
注意，从抽象角度看，是把一个新值赋给整个数组。

数组类型通常被视为无序的；但某些情况下，尤其是字符数组，我们希望定义一种与通常字典序一致的顺序。只有定义域类型和值域类型都有序时，才可能这样做。这时，两个数组的次序，由它们最小的不同下标处的元素次序决定。因此：

\quad ``BACK'' $<$ ``BANK''

\noindent
因为第三个字母是两者最先不同的地方，而且

\quad ``$C$'' $<$ ``$N$''

指定数组值的一种方便方法，是使用仿照 for 语句构造的\textit{for 表达式}：

\quad \textbf{for} $i:D$ \textbf{take} $E$

\noindent
其中，$E$ 是包含变量 $i$、产生值域类型值的表达式。当 $i$ 遍历定义域类型 $D$ 时，对 $E$ 求值，就得到相应数组元素的值。

如果数组的值域类型定义了某些操作，自然可以把这些操作推广到数组类型。例如，如果 $A$ 和 $B$ 是具有相同定义域的实数数组，就可以自然地写：

\quad $A + B,$ $A - B$,

表示定义域相同的数组，其元素分别是 $A$、$B$ 对应元素值的和与差。但程序员必须始终意识到，数组很大时，这些操作的成本可能很高；在从抽象程序推进到更具体程序的过程中，应当设法消除这些操作。

\subsection{表示}
\label{sec:representation-6.2}

大多数程序员都熟悉数组在计算机存储器中的表示。最常见的是非紧缩表示，为每个数组元素分配一个或多个完整的字。这种情况下，各元素的计算机地址可以简单地计算出来：先把下标值转换为最小表示，再乘以每个元素占用的字数，最后把结果加到数组第一个元素的地址上。利用计算机通常的按字选择机制，就可以访问和更新这个值，而不影响其他数组元素。

另一种表示把元素紧缩在字边界之内，使每个元素只占一个字内固定数量的位，虽然整个数组可能跨越多个字。以 monthtable 为例，每个元素只有四种取值，从 28 到 31；因此，在最小表示中只需两位就能容纳，整个数组也就可以容纳在二十四个连续的位中。

数组按这种方式紧缩时，选择带下标变量的值就复杂得多。为了选出正确的字，必须把下标的最小表示除以每个字包含的元素个数。把商加到数组第一个字的地址上，随后访问该字。把余数乘以每个元素的位数，用结果作为 $\text{shift-count}_2$，把所需值移到字内的标准位置。然后，可以用掩码去掉相邻数组元素中不需要的值。对紧缩数组元素进行选择性更新更加费事，因为必须把新值插到字内的正确位置，同时不改变相邻元素的值。如果每个字的元素个数恰好是二的幂，两种操作的效率都可以略有提高：下标的整数除法，可以用移位求商、用掩码求余数来替代。在某些机器上，把每个元素存储在单个字符位置，还能进一步提高效率。

数组的最小表示类似于笛卡尔积的最小表示，不过，每个元素值所乘的因子，等于值域类型基数的下标值次幂。除非值域类型的基数恰好是二的幂，否则选择或更新以最小表示存储的数组元素，比上面描述的过程还要费事。如果数组跨越多个普通计算机字，其代价就高得难以接受。因此，数组的最小表示主要只有理论意义。

\begin{figure}[h]
	\centering
	\begin{subfigure}[b]{0.3\textwidth}
		\centering
		\includegraphics[height=10em]{chap2/fig3a}
		\caption{}
	\end{subfigure}	
	\begin{subfigure}[b]{0.3\textwidth}
		\centering
		\includegraphics[height=8em]{chap2/fig3b}
		\caption{}
	\end{subfigure}	
	\begin{subfigure}[b]{0.3\textwidth}
		\centering
		\includegraphics[height=5.5em]{chap2/fig3c}
		\caption{}
	\end{subfigure}%
	\caption{$A:\text{\textbf{array} } 0\,.\,.\,7 \text{ \textbf{of} } T$ 的表示}
\end{figure}

有限映射的定义域本身是数据结构时，例如笛卡尔积，通常用最小表示来表示定义域，以免分配用不到的存储空间。例如，display page 的定义域，是整数范围 1 到 40 与 1 到 27 的笛卡尔积。采用最小表示，就得到从 0 到 $40 \times 27 - 1 = 1079$ 的整数范围。因此，分配 1080 个连续的字来保存数组元素的值。要访问给定行和字符位置处的元素，首先必须按第 \ref{sec:representation-4-2} 节介绍的方法，为下标构造最小表示。

多维数组的另一种表示方法，有时称为代码字方法或描述符方法，不过我们把它称为“树形表示”。其要点是分配一个一维\textit{基}数组，每个元素对应原数组的一行，并在其中存放一段连续存储区域的地址，该区域保存那一行的值。各行之间不必连续。这样，访问或更新元素时，就不必计算下标的最小表示。只需把行号加到树的基数组第一个元素的地址上，就能访问所需行第一个元素的地址；再把下一个下标值加到这个地址上，就得到所需元素的地址。

\begin{figure}[h]
	\centering
	\begin{subfigure}[b]{0.35\textwidth}
		\centering
		\includegraphics[height=12em]{chap2/fig4a}
		\caption{}
	\end{subfigure}	
	\begin{subfigure}[b]{0.6\textwidth}
		\centering
		\includegraphics[height=10.5em]{chap2/fig4b}
		\caption{}
	\end{subfigure}	
	\caption{二维数组的表示}
\end{figure}

在数组的非紧缩表示与紧缩表示之间选择，依据与笛卡尔积类似。如果需要快速访问和更新，就采用非紧缩表示；如果值域类型本来就适合放在计算机字边界内，例如元素是浮点数，这也是显然合适的选择。如果元素远短于一个字，并且存储空间不足，或数组的复制、比较比下标访问和选择性更新更频繁，就建议采用紧缩表示。紧缩数组一个特别常见的用途，是表示程序设计语言中的标识符。为了效率，可以把长到标准数组放不下的标识符截断，把过短的标识符用空白补齐。

多维数组表示方式的选择，依据则很不一样。标准表示更加节省存储空间，按行、按列或同时按行列遍历数组元素时，效率都很好。此外，如果数组必须整体在主存和辅助存储器之间传送，它也更方便。不过，在乘法较慢的机器上，采用树形表示更快，代价是接受保存地址数组所需的额外空间；只要每行不太短，这点空间并不大。如果每行只有两个字，数据存储就会多出百分之五十的开销。

选择树形表示还可能有几个其他理由：

\begin{enumerate}[wide, nosep, label=(\arabic*)]
	\item 在某些普遍采用动态存储分配的计算环境中，可能很难获得大块连续区域。这时，可以把大型二维数组拆成若干较短的行，方便地安排存储。

	\item 可以建立一种机制：处理某些行时，把其他行保存在辅助存储器中；当某一行不在主存中时，基数组中该行的主存地址由其辅助存储地址替代。这样，有望处理无法与处理程序一起完整放入主存的大数组。不过，必须仔细考察这种操作的成本，确保辅助存储传送次数可以接受。

	\item 在某些应用中，已知几个矩阵共享相同的行。树形表示可以只保存这种共享行的一份副本，复制它的地址，而不复制完整的值。但这种情况下，不能对共享行进行选择性更新。

	\item 即使对于一维数组，如果各元素的大小变化很大，也建议采用树形表示；对于多维数组，如果各行长度变化很大，同样建议采用它。
\end{enumerate}

\noindent
\textit{练习}
\nopagebreak

\noindent
某输入设备的字符集只有三十个字符，用枚举定义，其中包括 space、newline 和 newpage。可以调用过程

\quad read next character.

\noindent
从输入设备一次读入一个字符，放入缓冲区。应当逐行把这些字符组合到数组

\quad page:display page,

\noindent
中；收到 newpage 字符时，用指令

\quad outpage.

\noindent
把它输出到显示设备。显示设备\textit{不}识别 newline 或 newpage 字符，因此必须用空格填满行尾和页尾的剩余部分。

用合适的语言编写程序完成这一操作，并为 display page 选用不同表示，例如：

\quad 非紧缩表示

\quad 较松的紧缩表示

\quad 紧密的紧缩表示

\quad 间接表示。

\noindent
使用不同表示重写程序。比较不同表示涉及的代码和数据长度，以及运行速度。

编写相应程序，从显示设备读取一页，输出各个字符；注意尽可能去掉每行末尾多余的空格，以及每页末尾的空白行。
\section[幂集]{幂集}

给定集合的幂集，定义为该集合所有子集组成的集合；幂集类型的值则是集合，其元素从另一个类型中选取，这个类型称为幂集的\textit{基}类型。例如，原色已通过枚举定义为 red、yellow 和 blue。其他主要颜色由两种或三种原色混合而成：橙色是红色与黄色的混合，棕色是三种原色的混合。因此，每种主要颜色（包括原色）都可以用能混合出它的原色子集来指定。例如，橙色可以看作只有 red 和 yellow 两个元素的集合。用枚举定义集合的传统记法，可以写成 $\{$red, yellow$\}$。纯红色可以看作只有原色 red 一个元素的集合，即 $\{$red$\}$。这样，就能表示七种主要颜色：红、橙、黄、绿、蓝、紫和棕。没有任何原色时（即空集），可以把它看作表示没有颜色，例如白色。以这些颜色为取值的类型，可以声明为 primary color 类型的幂集：

\quad \textbf{type} color $=$ \textbf{powerset} primary color.

第二个例子来自表示电梯请求按钮状态所需的数据结构。floor 类型的简单变量（见第 \ref{sec:unstructured-data-types} 节），能够指出电梯的某一个停靠楼层。但如果要记录电梯内部整块按钮面板的状态，就需要用楼内所有可能楼层的一个子集来表示，即请求按钮已被按下的那些楼层组成的子集。因此，可以把 liftcall 类型定义为 floor 类型的幂集：

\quad \textbf{type} liftcall $=$ \textbf{powerset} floor.

第三个例子来自某种纸牌游戏，例如扑克或桥牌中的一手牌。一手牌是纸牌集合中没有重复元素的子集，因此适合用以下幂集类型中的一个值表示：

\quad \textbf{type} hand $=$ \textbf{powerset} cardface;

\noindent
这个类型涵盖至多五十二张牌的所有手牌，即使某种具体游戏会规定标准手牌张数，或把上限设为小于五十二的数。

最后一个例子表示计算机外部设备的状态，例如纸带读入机。尝试输入字符时，可能出现若干异常情况：

\begin{enumerate}[wide, nosep, label=(\arabic*)]
	\item 操作员把设备切换到“手动”。
	\item 没有装入纸带。
	\item 最近读入的字符发生奇偶校验错误。
	\item 最近读入字符时检测到纸带偏斜。
\end{enumerate}

\noindent
这些情况可以定义为一个枚举类型：

\quad \textbf{type} exception $= ($manual, unloaded, parity, skew$)$; 

\noindent
由于可能同时检测到其中几种情况，读入机的状态可以指定为某个幂集类型的值：

\quad \textbf{type} statusword $=$ \textbf{powerset} exception.

幂集类型的基数等于二的若干次幂，指数为基类型的基数，即：

\quad cardinality$(\text{\textbf{powerset} } D) = 2^{\text{cardinality} (D)}$

可以通过考察完整指定该类型的一个值需要作多少次决定，来证明这一点。对于基类型的每个值，都有两种选择：它在集合中，或不在集合中。这样的决定可以独立地作 cardinality$(D)$ 次。

\subsection{操纵}

集合的基本构造操作，是从定义域类型中取得若干值，把它们转换为一个恰好以这些值为元素的集合。与笛卡尔积一样，类型名用作转换函数；但对于集合，参数个数可以变化，从零个开始。例如：

\quad primary color$($red, yellow$)$ \tabto*{14em} 即橙色

\quad liftcall$($ground$)$ \tabto*{14em} 即只按下了一个按钮

\quad statusword$()$ \tabto*{14em} 即没有异常情况。

\noindent
后两个例子分别说明了\textit{单元素集合}（必须与其唯一元素明确区分）和不含任何元素的\textit{空集}的概念。如果在这种构造中省略类型名，就应按通常做法使用花括号，而不是圆括号。

与空集相对的是全集，它包含基类型中的所有值。可以记为：

\quad $T$.all.

\noindent
不过，只有基类型有限时，这个全集才作为可存储的数据值存在。

数学家和逻辑学家对集合的基本操作都很熟悉。

\begin{enumerate}[wide, nosep, label=(\arabic*)]
	\item 元素归属测试：如果 $x$ 在集合 $s$ 中，布尔表达式“$x$ \textbf{in} $s$”的值为 true，否则为 false。
	\item 相等：两个集合相等，当且仅当它们具有相同的元素。
	\item 交：$s_1 \wedge s_2$ 恰好包含同时在 $s_1$ 和 $s_2$ 中的那些值。
	\item 并：$s_1 \vee s_2$ 恰好包含在 $s_1$ 或 $s_2$ 中，或同时在两者中的那些值。
	\item 相对补集：$s_1 - s_2$ 恰好包含 $s_1$ 中不在 $s_2$ 里的那些元素。
	\item 包含关系测试：只要 $s_1$ 的所有元素也是 $s_2$ 的元素，$s_1 \subset s_2$ 的值就为 true，否则为 false。
	\item 集合的大小说明它有多少个元素。

	如果一个集合的定义域类型定义了某些运算符，构造相应的集合操作往往很有用。特别是，如果定义域类型有序，就可以使用以下运算符：

	\item $\min(s)$ 是 $s$ 的最小元素；$s$ 为空时未定义。
	\item $s \text{ \textbf{down} } n$ 是一个集合，恰好包含那些第 n 个后继在 $s$ 中的值。
	\item $s \text{ \textbf{up} } n$ 是一个集合，恰好包含那些第 $n$ 个前驱在 $s$ 中的值。
	\item Range$(a,\ b)$ 在 $a \leqslant b$ 时是包含 a、succ$(a), \dots, b$ 的集合，否则为空。
\end{enumerate}

集合最有用的选择性更新操作是：

\noindent
$x:\vee\ y$; \tabto*{5em} 把集合 $y$ 并入 $x$

\noindent
$x:\vee\ T(a)$ \tabto*{5em} 把元素 $a$ 加入 $x$

\noindent
$x:\wedge\ y$; \tabto*{5em} 从 $x$ 中删除所有不属于 $y$ 的元素

\noindent
$x:-\ y$ \tabto*{5em} 从 $x$ 中删除所有也属于 $y$ 的元素

\noindent
\hangindent=3.2\parindent$x:\text{\textbf{down} }n$ \tabto*{5em} 把 $x$ 的每个元素减去 $n$，不能这样做的元素则删除

\noindent
\hangindent=3.2\parindent$x:\text{\textbf{up} } n$ \tabto*{5em} 把 $x$ 的每个元素加上 $n$，不能这样做的元素则删除

\noindent
有时，从 $x$ 中选出一个元素，并同时把它从 $x$ 中移除，也很有用。这个操作可以记为：

\quad $a$ \textbf{from} $x$.

\noindent
如果 $x$ 的定义域类型有序，自然应选取 $x$ 的最小元素；否则应把选择视为任意的。

常常可以用某个条件 $B$ 来定义集合的值：定义域类型中恰好满足 $B$ 的那些值，就是所需集合的元素。可以记为：

\quad $\{i:D \mid B\}$

\noindent
其中 \tabto*{4em} $i$ 是类型 $D$ 的变量，视为 $B$ 的局部变量，

\noindent
而 \tabto*{4em} $B$ 是一个布尔表达式，通常包含并依赖于 $i$。

\noindent
要使这个表达式表示幂集类型的一个值，$D$ 的基数必须有限，而且 $B$ 必须对该类型的所有值都有定义。

最后，经常需要对某个集合的每个元素进行某种操作，也就是执行一个循环，让计数变量依次取到集合中的所有值。合适的表达记法是：

\quad \textbf{for} $x$ \textbf{in} $s$ \textbf{do} \dots

\noindent
如果 $s$ 的基类型有序，规定从最小元素开始，按自然顺序选取各元素，是合理的。对于无序的基类型，程序员不关心按什么顺序选取元素，因此保留选择最有利于效率的顺序的余地。

\subsection{表示}

为幂集选择计算机表示时，希望所有基本操作都能简单地用单条机器码指令执行，同时尽量减少存储占用。对于大多数数据结构存储方法，这两个目标存在根本冲突，因此程序员必须在不同表示方法之间作出选择；但对于幂集，只要基类型不太大，两个目标就能完全兼顾。

建议的表示方法，是在存储器中分配与集合潜在元素个数相同数量的位。基类型的每个值都对应一个位；如果它确实属于集合，这个位为一，否则为零。例如，color 类型的每个值都能用三位表示；最高有效位对应原色 red，最低有效位对应 blue。因此，橙色表示为 110，红色表示为 100。大小为 $n$ 的每个集合，都表示为一个恰好在对应位置上有 $n$ 个一的位模式。因此，空集用全零位模式表示。

另一个例子是 hand 类型，它的表示需要五十二位，每一位对应 cardface 类型的一个值。这里，最好对基类型采用最小表示，以免位模式表示中留下不用的空隙。

由于幂集类型的值的个数总是二的整数次幂，对于基类型较小的幂集，在二进制计算机上不可能有比位模式表示更节省存储空间的方法。接下来还要说明，幂集类型上的操作也能高效执行。

\begin{enumerate}[wide, nosep, label=(\arabic*)]
	\item 要得到以 $x$ 为唯一元素的集合，可以在符号位位置装入一个 1，再向右移动 $x$ 位。在移位较慢的计算机上，可以用查表达到同样效果。用分量构造集合，可以把相应的所有单元素集合作逻辑并。

	\item 元素归属测试 $x$ \textbf{in} $s$，可以把 $s$ 向高位方向移动 $x$ 位，再检查最高有效位：1 代表 \textbf{true}，0 代表 \textbf{false}。

	\item 二进制计算机常常提供单条指令，执行逻辑交、并和取补。

	\item 有时可以用内置的机器码指令，统计一个字中值为 1 的位数，从而得到集合的大小。否则，可以反复执行规格化，每次用掩码清除紧邻符号位的那个位，来确定集合的大小。第三种办法是把位模式拆成若干小块，对每块查表，再把结果相加。

	\item \textbf{up} 和 \textbf{down} 操作显然可以用右移或左移完成。

	\item 可以用规格化指令高效找到集合的 min；这条指令自动计算，把第一个为一的位移到紧邻符号位的位置，需要移动多少位。

	\item for 语句也可以利用规格化高效实现，每遇到一个为一的位，就用掩码把它清除。

	\item range 操作可以用两次移位完成，第一次移位会重新补入符号位。	
\end{enumerate}

因此，只要定义域类型的基数，不大于计算机中可施加逻辑和移位操作的最大字的位数，所有这些操作都能高效执行。如果占用空间明显超过一个这样的字，使用选择性更新操作通常比使用一般的、产生结果值的运算符更划算。此外，size 和 min 这类操作可能变得相当低效。此时，往往可以把这些值与集合一起冗余保存，每当集合更新就同步维护，而不在每次需要时重新计算。

如果已知基类型的基数很大（甚至可能无限），而典型集合的大小相对很小，那么位模式表示就完全失去了吸引力，因为存储和操作大片零已经不划算。这种稀疏集合的处理留到第 \ref{sec:sparse-data-structures} 节再讨论。

\subsection{例子}

\noindent
\textit{问题}：编写程序，构造集合

\quad primes:\textbf{powerset} $2\,.\,.\,N$;

\noindent
使其包含基类型中的所有素数。

\noindent
采用埃拉托色尼筛法，避免所有乘法和除法。

埃拉托色尼的方法，是先把所有数放入“筛子”，然后重复以下步骤，直到筛子为空：

\noindent
选出并移除筛子中剩下的最小数（它必然是素数），然后逐步遍历筛子，移除这个数的所有倍数。

程序很容易写出：

\noindent
\textbf{begin} $n$, next:$2\,.\,.\,N$; sieve:\textbf{powerset} $2\,.\,.\,N$;

\noindent
\quad sieve $\coloneq$ range$(2, N)$;

\noindent
\quad primes $\coloneq \{\}$;

\noindent
\quad \textbf{while} sieve $\neq$ empty \textbf{do}

\noindent
\quad \quad \textbf{begin} next $\coloneq$ min$($sieve$)$;

\noindent
\quad \quad \quad primes: $\vee\ \{$next$\}$;

\noindent
\quad \quad \quad \textbf{for} $n\coloneq$ next \textbf{step} next \textbf{until} $N$ \textbf{do}

\noindent
\quad \quad \quad \quad sieve:$-\ \{n\}$

\noindent
\quad \quad \textbf{end}

\noindent
\textbf{end} primefinder.

但如果 $N$ 很大，例如在 $10\,000$ 的数量级，这个程序直接执行时就无法达到可以接受的效率。解决办法是把它作为算法的抽象模型，改写为更高效的形式，只使用不超过计算机字长的集合上的操作。因此，需要声明一个字数组来表示这两个集合，并假定“wordlength”是一个环境查询值，给出一个字包含的位数：

\quad primes, sieve:\textbf{array} $0\,.\,.\,W$ \textbf{of powerset} $0\,.\,.\,\text{wordlength} - 1$

\noindent
其中 $W = (N + 1) \div \text{wordlength} + 1$。

\noindent
这意味着两个集合可能略大于 $N$，但为方便起见，接受这种无害的扩展。

为了访问这些集合中的一个位，必须同时知道字号和位号。由于不想用除法来求它们，我们把计数变量 $n$ 和 next 表示为笛卡尔积：

\quad $n$, next:$(w,\ b$:integer$)$;

\noindent
其中 $w$ 表示字号，$b$ 表示位号。

现在不妨先重新编码最内层循环，检查这种表示的效率。

\quad \textbf{for} $n\coloneq$ next \textbf{step} next \textbf{until} $N$ \textbf{do} sieve:$-\ \{n\}$;

\noindent
改写为：

\quad $n\coloneq$ next;

\quad \textbf{while} $n.w \leqslant W$ \textbf{do}

\quad \quad \textbf{begin} sieve$[n.w]$:$-\ \{n.b\}$;

\quad \quad \quad $n.b\coloneq n.b + \text{next}.b$;

\quad \quad \quad $n.w\coloneq n.w + \text{next}.w;$

\quad \quad \quad \textbf{if} $n.b \geqslant$ wordlength \textbf{then begin} $n.w \coloneq n.w + 1$;

\quad \quad \quad \quad $n.b \coloneq n.b - \text{wordlength}$

\quad \quad \quad \textbf{end}

\quad \quad \textbf{end}

\noindent
看来这样的效率可以接受，于是继续编码外层循环的其他操作，从最困难的开始：

\quad next $\coloneq$ min$($sieve$)$;

\noindent
这里，不希望每次都从 sieve 集合的起点开始搜索最小值，因为过程接近结束时，这样会扫描许多空字。因此，我们利用这一事实：next 的新值必定大于旧值。

搜索分两步：先找一个非空字，再找其中第一个为一的位。但如果搜索字时到达数组末尾，整个程序就完成了。

\quad \textbf{while} sieve$[\text{next}.w] = \{\}$ \textbf{do} $\{\text{next}.w \coloneq \text{next}.w +1$;

\quad \quad \textbf{if} next.$w > W$ \textbf{then exit} primefinder$\}$;

\quad \quad next.$b\coloneq$ min$(\text{sieve}[\text{next}.w])$;

其他操作都很简单。由于外层循环用 exit 终止，就不必再单独测试 while 条件；语句

\quad primes:$\vee\ \{$next$\}$;

\noindent
可以编码为：

\quad primes$[\text{next}.w]$:$\vee\ \{\text{next}.b\}$.

包含初始化的整个程序如下：

\quad primes, sieve: \textbf{array} $0\,.\,.\,W$ \textbf{of powerset} $0\,.\,.\,\text{wordlength} - 1$;

\quad \textbf{begin} primefinder;

\quad \quad $n$, next:$(w, b$:integer$)$;

\quad \quad \textbf{for} $t$:$0\,.\,.\,W$ \textbf{do begin} primes$[t] \coloneq \{\}$;

\quad \quad \quad sieve$[t] \coloneq$ range $(0\,.\,.\, \text{wordlength} - 1)$

\quad \quad \textbf{end};

\quad \quad sieve$[0]$:$-\ \{0, 1\}$;

\quad \quad next.$w\coloneq 0$;

\quad \quad \textbf{while} true \textbf{do}

\quad \quad \quad \textbf{begin while} sieve$[\text{next}.w] = \{\}$ \textbf{do}

\quad \quad \quad \quad \textbf{begin} next.$w\coloneq$ next.$w + 1$;

\quad \quad \quad \quad \quad \textbf{if} next.$w > W$ \textbf{then exit} primefinder

\quad \quad \quad \quad \textbf{end};

\quad \quad \quad \quad next.$b\coloneq$ min$($sieve$[\text{next}.w])$;

\quad \quad \quad \quad primes$[\text{next}.w]$: $\vee\ \{\text{next}.b\}$;

\quad \quad \quad \quad $n\coloneq$ next;

\quad \quad \quad \quad \textbf{while} $n.w \leqslant W$ \textbf{do}

\quad \quad \quad \quad \quad \textbf{begin} sieve$[n.w]$:$-\ \{n.b\}$;

\quad \quad \quad \quad \quad \quad $n.b\coloneq n.b + \text{next}.b$;

\quad \quad \quad \quad \quad \quad $n.w\coloneq n. w + \text{next}.w$;

\quad \quad \quad \quad \quad \quad \textbf{if} $n.b \geqslant$ wordlength \textbf{then}

\quad \quad \quad \quad \quad \quad \textbf{begin} $n.w\coloneq n.w + 1$;

\quad \quad \quad \quad \quad \quad \quad $n.b\coloneq n.b - \text{wordlength}$

\quad \quad \quad \quad \quad \quad \textbf{end}

\quad \quad \quad \quad \quad \textbf{end}

\quad \quad \quad \textbf{end}

\quad \textbf{end} primefinder

这个程序的一个特点，是使用环境查询 word length，充分发挥机器所能达到的效率，同时又完全独立于机器。即使在字长差别很大的机器上，程序也不仅能运行，而且能高效运行。

但这个程序最值得关注的特点，是它与前一版本的关系。从抽象角度看，两者表达的是同一个算法；改变的只是计算机上的数据表示方式。最初的设计充当了框架或范式，第二个版本更复杂的编码就是据此组织的。分两个阶段进行设计，简化了这样一项任务：确保最终程序的每个部分都能与其他部分成功配合。

\noindent
\textit{练习}
\nopagebreak

\noindent
使用只表示奇数的集合重写程序。（提示：先重写更抽象的程序。）
\section[序列]{序列}

前面的章节讨论了初等数据结构。它们在实际程序设计中非常重要，在现代数字计算机上表示和操纵时，也几乎不造成困难。此外，它们为构造一切更高级的结构提供了不可缺少的基础。

初等结构化类型与高级结构类型最重要的区别在于：只要组成类型的基数有限，前者的基数就严格有限。有限集合与无限集合的区别，具有深刻的数学意义，也给表示和操纵方法带来许多影响。

\begin{enumerate}[wide, nosep, label=(\arabic*)]
	\item 高级结构类型的潜在值可能有无限多个，因此，保存某个值所需分配的存储量，无法仅凭声明来确定。通常只有程序实际运行时才能确定，而且在许多情况下，它会在执行过程中变化。对于初等结构，不同潜在值的个数有限，保存任意值所需的最大存储量是固定的，可以由声明形式确定。

	\item 结构化值的大小相当大时，分别更新结构的各个分量，比给整个结构赋一个全新值更高效。即使对于初等类型，有时也发现选择性更新更高效，尤其是非紧缩表示的笛卡尔积和数组。对于高级数据结构，选择性更新的效率优势通常更加明显。

	\item 大小动态变化的高级数据结构，需要某种动态存储分配和释放机制。所需的存储单元通常用指针连接起来，指针有时也称为引用或地址；释放则通过显式编程操作，或某种通用垃圾回收来完成。动态存储分配和指针的使用，显著增加了处理复杂性；当数据必须在计算机主存与辅助存储器之间传送时，问题尤其严重。对于初等数据结构，这类问题完全可以不出现。

	\item 为高级数据结构选择合适表示，往往比为初等结构选择表示困难得多。各种基本操作的效率强烈依赖表示选择，因此，要作出合理选择，就需要知道这些操作的相对调用频率。结构的一部分或全部保存在辅助存储器上时，这一点尤其重要；此时，表示的选择还应考虑硬件设备的特性，即旋转介质上的磁道和柱面排列，以及磁头移动时间和旋转延迟。对于初等结构，在大多数表示中，各基本操作的效率大致相当。
\end{enumerate}

由此可见，高级结构与初等结构差别十分明显，高级结构带来的问题也大得多。这说明，对于实际程序员，明智的做法是尽可能只使用初等结构，只有应用本身迫使他这样做时，才使用高级结构。

高级数据结构中，第一个也是最熟悉的例子是序列。它仅仅被看作某个给定类型的任意多个项按顺序组成的结构。“序列”一词意在涵盖磁带、磁盘、磁鼓或主存中的序列。主存中的序列，有时称为流、表、串、栈、双端队列、队列，甚至集合。文件（或顺序文件）一词，常用于保存在辅助存储器中的序列。序列概念是一种抽象，所有这些结构都可以视为它的不同表示。

序列的第一个例子是串，ALGOL 和 SNOBOL 程序员对它很熟悉。串由任意长度的字符序列构成，因此可以定义为：

\quad \textbf{type} string $=$ \textbf{sequence} character.

第二个例子来自数据处理应用，即维护一个汽车数据文件。文件中的每个项（有时称为记录）表示一辆汽车，因此属于 car 类型；前面已经给出了 car 类型的一种可能定义：

\quad \textbf{type} car file $=$ \textbf{sequence} car.

第三个例子给出处理一副牌的另一种方法。它可以看作一个纸牌序列，其长度可能随发牌而变化：

\quad \textbf{type} deck $=$ \textbf{sequence} cardface;

\noindent
当然，并非所有纸牌序列都表示现实中的牌组。例如，同一张牌出现两次的序列无效，程序员应当避免。因此，有效牌组的最大长度是 52，虽然声明中没有表达这一事实。

下一个例子来自一类特定符号表达式的处理，即多项式。多项式

$a_n x^n + a_{n-1}x^{n-1}\dots a_1 x + a_0$

\noindent
可以表示为其系数 $a_i$ 的序列。如果多项式次数 $n$ 无法预先确定，或在计算过程中变化，那么序列就是最合适的定义方式：

\quad \textbf{type} polynomial $=$ \textbf{sequence} integer.

最后一个例子说明，如何表示程序设计语言中的标识符概念。理论上，标识符可以有任意长度，因此需要序列。序列中的项要么是字母，要么是数字。不过，首字符总是字母，可以与其余部分分开。因此，精确定义与标识符对应的数据结构，可以写成：

\quad \textbf{type} identifier $= ($first:letter; rest:\textbf{sequence} $(l$:letter, $d$:digit$))$.

\subsection{操纵}

序列类型 $T$ 的零元素是不含任何项的序列\textemdash{}称为空序列，记为 $T()$。对于定义域类型的每个值 $v$，都有一个仅含 $v$ 这一项的序列，称为 $v$ 的\textit{单项序列}，记为 $T(v)$。最后，如果 $v_1, v_2, \dots, v_n$ 是基类型的值（可能重复），
$T(v_1, v_2, \dots, v_n)$ 就表示由这些值按给定顺序组成的序列。为方便起见省略类型名 $T$ 时，我们用方括号把序列括起来：

\quad $[v],$\qquad $[v_1, v_2, \dots, v_n]$

\noindent
不过，字符序列通常用引号括起来表示。

序列的基本操作是串接，即把两个序列一前一后接在一起。因此，如果 $x$ 是字符序列“PARIS IN THE”，$y$ 是序列“THE SPRING”，它们的串接 $x\wideparen{\ \ }y$ 就是序列：

\quad $z =$ ``PARIS IN THETHE SPRING''

\noindent
除非操作数非常小，否则计算机上的串接效率很低，因为通常需要为两个操作数都创建新的副本。因此，程序员应尽力用选择性更新代替串接。

把序列分解为组成部分的基本运算符，有返回非空序列第一项和最后一项的：

\quad $x$.first, $x$.last

\noindent
也有移除非空序列的最后一项或第一项，分别产生前段或后段的：

\quad initial$(x)$, final$(x)$.

序列之间的一种重要关系，是序列 $x$ 等于序列 $y$ 的某个起始子序列或末尾子序列：

\quad $x$ \textbf{begins} $y$

\quad 或 $x$ \textbf{ends} $y$。

\noindent
在前面的例子中，“PARIS” \textbf{begins} $z$，“RING” \textbf{ends} $z$。这两种测试在程序运行中可能相当费时，应尽可能避免。

序列的一个重要性质是长度，即它包含的项数；对于序列 $x$，可以用函数 length$(x)$ 求得。

对于某些用途（例如构造词典），按传统字典序原则把序列类型视为有序的，会很有用。与数组一样，两个序列的次序，由它们第一处不同项的次序决定；如果没有这样的项，则较短的序列排在以它为起始部分的较长序列之前，例如：

\quad ``ALPHA'' $<$ ``ALPHABET''。

\noindent
在这种顺序中，每个序列都有后继，但只有少数序列有前驱。

序列上一种极重要的选择性更新操作，是在现有序列 $x$ 的末尾追加一个新值 $v$。可以写成：


\quad $x\wideparen{\text{: \ } }T(v)$;

\noindent
它对应于把值 $v$ 写入顺序文件 $x$ 这一熟悉概念。与读取文件 $x$ 开头相对应的操作，则是移除 $x$ 的第一项，并把它的值赋给某个变量 $v$。可以写成：

\quad $v$ \textbf{from} $x$;

\noindent
某些应用中，能够从序列中读回最近写入的项会很有用；可以表示为：

\quad $v$ \textbf{back from} $x$;

\noindent
它从 $x$ 中移除最后一项。普通写入操作

\quad $x\wideparen{\text{: } T}(v).$

\noindent
把项“压入”栈，而这个读回操作则可以把栈顶项“弹出”。需要时，还可以定义第四种更新操作，即把新值添加到序列开头：$($putback$(x$, $v))$。

某些情况下，避免 \textbf{from} 操作中涉及的项复制，效率会更高。只需省略左边的变量，就可以处理这些情况，例如：

\quad \textbf{from} $x$

\quad \textbf{back from} $x$。

\noindent
此时，通常通过选择函数 $x$.first 和$/$或 $x$.last 访问序列中的项。

依次扫描序列中的所有项，是很常见的需求；可以仿照 for 语句采用如下记法：

\quad \textbf{for} $v$ \textbf{in} $x$ \textbf{do} $S$;

\noindent
如果 $x$ 为空，就不执行该语句。否则，变量 $v$（视为 $S$ 的局部变量）依次取序列 $x$ 中各项的值，每取一个值就执行一次 $S$。在这一构造中，$S$ 内部不应更新 $x$ 或 $v$。

类似构造可以用来定义一个序列：对序列中的每一项 $v$ 逐项施加变换 $E(v)$。

\quad \textbf{for} $v$ \textbf{in} $s$ \textbf{take} $E(v)$。

决定序列的表示时，最重要的是知道将对它进行哪些选择性更新操作。

\begin{enumerate}[wide, nosep, label=(\arabic*)]
	\item 如果唯一的操作是 \textbf{from}，这个序列称为\textit{输入}序列。显然，要让输入序列具有任何值，都必须把它初始化为声明它的外部环境中的某个现有值。把程序的局部序列，与辅助存储器上比较持久存在的某个文件关联起来，通常称为为输入而“打开”文件。我们假定，声明局部输入序列时会隐式调用这个操作。退出该序列所属的程序块时，则隐式调用相反的“关闭”文件操作。

	\item 如果唯一的操作是向文件写入，这个序列称为\textit{输出}序列。输出序列可以像输入序列一样，从环境中取得初值；更常见的是，初值为空。不论哪种情况，要让它有实际用途，都必须在退出程序块时，把序列的最终值赋给声明它的外部环境中的某个变量。声明序列时，应同时指出这个外部变量；如果该变量比较持久地保存在辅助存储器中，就称为输出文件。在进入和退出程序块时，隐式调用文件打开和关闭操作的规则，与输入文件类似。

	\item 如果只有写入和读回操作（压入和弹出），序列称为\textit{栈}；栈的初值总为空，最终值通常不保存。

	\item 如果只有在末尾写入、从开头读取这两种操作，序列称为\textit{队列}；同样，初值总为空，最终值通常不保存。

	\item 如果允许在序列两端读写，它有时称为\textit{双端队列}（deque）。不过，要让四种操作同样高效，表示就需要一定复杂性；幸运的是，大多数程序都可以不用双端队列。
\end{enumerate}

\subsection{表示}

\subsubsection{连续表示}

表示序列最简单的方法，是为它分配一块固定、连续的存储区域，足以容纳实际需要的所有项。如果序列的值（或至少其长度）在整个程序执行中保持不变，就适合这种方法\textemdash{}例如，用作输出消息或标题的字符串。

某些情况下，编写程序时还不知道序列长度，但在进入声明该序列的程序块时就已知道，而且这一长度在序列整个存在期间保持不变。此时，可以使用标准的栈式存储分配和回收方法，在程序块的局部工作区中分配一块连续存储区域。

即使序列长度会变化，有时也可以为它规定一个足够小、能够接受的上界，并永久分配这一最大存储区域。如果程序运行时超过上限，程序员必须愿意接受程序立即终止。除了固定区域，还需要指针或计数来指出序列当前的起点和终点。对于栈，第一项总在起点，因此只需要一个指向栈顶的指针。对于队列，序列有时会越过存储区域的末端，转回区域开头继续。这种表示称为循环缓冲区，可以用于并行程序设计，让并行运行的进程交换信息。此时，写进程发现缓冲区已满，就必须等待，直到读进程再次减小序列的大小。同样，缓冲区为空时，读进程必须等待。

连续表示最合适的另一种情况，是程序只需一个序列。因此，在为其他用途分配存储空间以后，它可以占用全部剩余空间；如果发生溢出，那么程序无论如何都不可能运行。如果需要两个栈，也可以容纳：让一个从剩余存储空间的一端开始，向上增长，另一个从另一端开始，向下增长。如果两个栈相遇，程序就无法继续。

如果要表示许多序列，可以建立一种机制，把它们分散放在剩余存储空间中；当某个序列增长到与邻居相遇时，重新移动安排部分或全部序列，让每个序列又有足够空间继续增长一段。每个序列都必须有一个指向其起点的\textit{基址单元}，并始终通过它来寻址。此外，还必须存储序列的实际长度。每次扩展序列时，都要检查相邻序列的基址单元和长度。重新安排时，要更新所有已移动序列的基址单元，指向序列的新位置。如果已知重新安排相对不频繁，这种专用机制相当有用；否则，应优先采用非连续表示。

\begin{figure}[h]
	\centering
	\begin{subfigure}[b]{0.15\textwidth}
		\centering
		\includegraphics[height=3.4em]{chap2/fig5a}
		\caption{}
	\end{subfigure}	
	\begin{subfigure}[b]{0.3\textwidth}
		\centering
		\includegraphics[height=8.1em]{chap2/fig5b}
		\caption{}
	\end{subfigure}	
	\begin{subfigure}[b]{0.45\textwidth}
		\centering
		\includegraphics[height=9em]{chap2/fig5c}
		\caption{}
	\end{subfigure}	
	\caption{序列（连续表示）}
\end{figure}

序列各项长度可变时，通常不必把较短项填充到最大长度，因为可以用判别标记字段或其他技术指出任一项的长度；读入该项时，便可以据此让指针移动正确的距离。但这要求在写入时就知道读取方向，如栈或队列的情况。如果要从两端读取，就必须保证，既能从项的开头推断其长度，也能从末尾推断；这需要在序列各项之间存储冗余信息（例如前一项的长度）。

\begin{figure}[h]
	\centering
	\begin{subfigure}[b]{0.3\textwidth}
		\centering
		\includegraphics[height=8.4em]{chap2/fig6a}
		\caption{}
	\end{subfigure}	
	\begin{subfigure}[b]{0.3\textwidth}
		\centering
		\includegraphics[height=8em]{chap2/fig6b}
		\caption{}
	\end{subfigure}	
	\begin{subfigure}[b]{0.3\textwidth}
		\centering
		\includegraphics[height=10em]{chap2/fig6c}
		\caption{}
	\end{subfigure}
	
	\begin{subfigure}[t]{0.15\textwidth}
		\centering
		\includegraphics[height=.4em]{chap2/fig6tag}
	\end{subfigure}
	\caption{序列（变长项）}
\end{figure}

一个序列本身是另一个序列某一项的一部分时，可以对这个项内序列采用连续表示。通常还要附一个计数，给出序列长度，以便读取项时计算其实际大小。

\subsubsection{链式表示}

为避免上一小节提到的重新移动安排问题，通常引入间接或链式存储分配方法，分配单位可以是固定长度，也可以是可变长度。可用存储空间被划分为若干区域，其中一些正用来保存某些序列的项，另一些空闲。空闲区域也链接成一个链式序列。每当程序员的序列需要扩展，就从空闲链中取得一个区域（或区域的一部分）；每当读取使序列缩短，就可以把一个区域归还空闲链。对于定长项，采用显式释放来管理动态存储分配并不困难。这里不讨论变长分配的问题；最好用分块来避免它们（见下一小节）。

最简单的链是单链。序列的每一项都附带一个链接单元，存放链中下一项的地址。空序列用一个不可能是地址的值表示（例如零或负一）；序列最后一项的链接单元也存放这个值。序列的基址单元指向链的第一项。

已知序列读取方向时，单链很有用，因为链接必须指向这个方向。对于栈，链接向后指；对于输入序列、输出序列和队列，则向前指。对于输入或输出序列，用于保存序列初值和$/$或终值的外部变量，其基址单元始终指向链的开头，而序列自身的基址单元则逐步沿序列前进。对于队列，使用两个基址单元，分别指向序列的两端。

\begin{figure}[h]
	\centering
	\begin{subfigure}[b]{0.3\textwidth}
		\centering
		\includegraphics[height=12em]{chap2/fig7a}
		\caption{}
	\end{subfigure}
	\begin{subfigure}[b]{0.3\textwidth}
		\centering
		\includegraphics[height=12em]{chap2/fig7b}
		\caption{}
	\end{subfigure}
	\begin{subfigure}[b]{0.3\textwidth}
		\centering
		\includegraphics[height=12em]{chap2/fig7c}
		\caption{}
	\end{subfigure}
	\caption{序列（链式表示）}
\end{figure}

对于栈，单链表示的一个可能优点，是多个栈可以共享相同的起始部分，从而在某些应用中节省空间和时间。不过，从这种栈中弹出一项时，不能立即把它占用的存储空间归还空闲链，因为它可能仍被其他栈使用。一种解决办法是从不显式归还空间，而是等空闲链耗尽。然后扫描所有当前已分配的序列，标记所有正在使用的块，再把所有未标记的块收集到空闲链。这称为扫描—标记—收集式垃圾回收。它看起来免除了程序员显式控制主存分配和释放的责任，但在复杂的计算机应用中，这种责任不能轻易回避，因此这样做可能有危险。

对于需要双向读取的双端队列，单链已经不够。通常的解决办法，是为链中的每一项附加\textit{两个}指针，一个指向前一项，另一个指向后一项。实际上，这两个指针可以压缩到一个字中，只存它们的差。由于第一项和最后一项中的一个指针为标准空值，总可以通过减法得到这些项中另一个指针的值。读写时，可以方便地调整新的第一项或最后一项的链接单元值，因为此时仍知道原来第一项或最后一项的地址。这一方案的具体推导留作练习。

链接各项的另一种方法，是把所有链接汇集到一个连续的表中，最好是定长表。这就给出了序列的一种树形表示，便于双向扫描。但它给序列的项数设定了上限，而且意味着链接所用的存储单元必须永久分配，即使序列很短也是如此。可以用分块减轻这一问题。

\subsubsection{分块表示}

链式表示的一个缺点，是需要额外空间保存链接，而且每次操作都要花时间管理空闲存储链。序列各项很小而序列很长时，这些问题尤其严重。解决方法是分块，即把连续存储和链式技术结合起来。

这种技术分配一个定长存储块，足以容纳大约十到一百个项。块填满后，用上一小节介绍的任意方法把新块链接到它后面。输入时，只有一个块中的所有项都被扫描完，才把该块归还空闲存储区。这样，链接占用的存储量可以降到微不足道的程度。对链的树形表示，这尤其有利。

如前所述，分块还能避免变长动态存储分配带来的问题，因为可以让所有序列都采用相同的固定块大小，不受各项大小的影响。此外，如果序列的一部分或全部保存在辅助存储器上，几乎都应当采用分块，因为传送单位太小时，辅助存储传送可能十分低效。分块唯一的（也值得怀疑的）缺点，是妨碍有效共享栈的尾部。

剩下的唯一问题，是选择适合各种用途的块大小。显然，它必须足够大，能容纳任一序列中的最大项。实际上，它还应大到至少容纳十个典型项；否则，块末尾不足以容纳下一项的剩余空间，可能占相当大的比例。此外，如果序列部分或全部保存在辅助存储器上，块应足够长，保证不会在块间间隙中浪费过多空间；传送频率也应足够低，保证不会把过多时间耗在启停、等待或磁头移动延迟上。

另一方面，如果块太大，第一个块开头和$/$或最后一个块末尾浪费的空间就会变得显著。因此，块也应足够小，保证典型序列至少占十个块。

面对这么多相互冲突的考虑，很难为长度和项大小各不相同的序列，以及采用不同访问方法的各种辅助存储器，选择统一的块大小。不过，往往可以作出一个可接受的折中；在当今计算机设计中，128 到 1024 个字的块通常是合适选择。大多数情况下，选定的大小即使增至两倍或减半，可能也不会有决定性影响。

\subsubsection{辅助存储表示}

处理序列时，程序通常只需访问它的一端，而中间和另一端的内容会在相当长的时间内不被使用。如果主存略有不足，把这些内容转移到辅助存储器就很有利，它们占用的主存空间可以用于其他用途。对于生存期长于读写它们的程序的输入和输出序列，几乎必须用辅助存储器长期保存。

采用辅助存储器时，处理和表示的效率要求按大小合理的块传送。包含序列活动端的块始终保留在主存中；为了让输入$/$输出与计算重叠进行，写入时的前一个块，或读取时的下一个块，在传送期间也保留已分配的空间。这称为双缓冲。还可以在存储器中保留更多缓冲区，平滑处理速度与传送速度的波动；但程序设计者不能误以为，这能解决处理速度与传送速度根本不匹配的问题。一般而言，如果双缓冲或三缓冲不够，再用更多缓冲区填满存储器，也不值得。

机器具有自动分页机制时，序列表示的问题会大大减轻。对程序员而言，只需采用连续表示，为最长可能序列分配所需空间。这实际上不应浪费存储，因为分页系统应推迟实际分配，直到空间第一次使用。随着序列扩展，会分配新的存储块，但在程序员看来，这些块的地址是连续的；因此，不会有块末空间不足以容纳下一项而闲置的问题。一个块填满后不久，就会自动迁移到辅助存储器；一旦需要，又会自动调回。对于输入序列，扫描完一个块以后不久，它也会从主存移走；但如果自上次读入以来内容未变，就不会产生不必要的辅助存储传送。分页系统唯一不会自动进行的操作，是在实际需要之前，预先把输入序列的一个块读入主存。
\section[递归数据结构]{递归数据结构}

数据结构化的方法，与处理这些数据的程序的结构化方法之间，存在一些密切的类比。因此，笛卡尔积对应复合语句，后者给各分量赋值。类似地，带判别标记的联合对应条件构造或 case 构造，为每种备选情况选择合适的处理方法。数组和幂集对应依次遍历其元素的 for 语句，迭代次数本质上有界。

序列结构是第一个允许构造基数无限、值的长度无界的类型的结构；它对应无界循环，用 \textbf{while} 条件控制终止。序列之所以无界，是因为构成它的一个分量（即前段）属于与自身相同的类型；正如 \textbf{while} 循环执行一次迭代以后，剩下要执行的语句与迭代前是同一条语句。

自然会问：这种类比能否推广到与递归过程相对应的数据结构？这种类型的一个值，可以包含不止一个与自身同类型的分量；正如递归过程可以在过程体中的多个位置递归调用自身。与递归过程一样，可以直接把正在定义的类型名写进其自身定义，方便地定义这种结构；对于相互递归的定义，也可以把类型名写在先前某个类型的定义中。

递归数据结构最明显的例子，来自对程序设计语言中算术表达式或逻辑表达式的描述；这里的递归反映了一个表达式可以嵌套在另一个表达式中的可能性。例如，算术表达式可以定义如下：

“表达式是一系列项，每一项由一个符号 $(+ \text{ 或 } -)$ 后接一个因子序列组成。除第一个因子外，每个因子都由一个符号 $(\times \text{ 或 } /)$ 后接一个基本表达式（primary）组成。基本表达式要么是常量，要么是变量，要么是括号括起来的表达式。表达式开头的正号可以省略。”

只需使用已经熟悉的技术，再加上递归，就能定义一个以这类表达式为值的结构化数据类型：

\quad \textbf{type} expression $=$ \textbf{sequence} term;

\quad \textbf{type} term $= ($addop:operator; $f$:\textbf{sequence} factor$)$;

\quad \textbf{type} factor $= ($mulop:operator; $p$:primary$)$;

\quad \textbf{type} primary $= ($const:$($val:real$)$,

\tabto*{9.3em}var:$(id$:identifier$)$,

\tabto*{9.3em}bracketed:$(e$:expression$))$;

\quad \textbf{type} operator $= ($plus, minus, times, div$)$;

这个定义表达算术表达式的抽象结构，却不涉及其作为字符串的具体表示细节。例如，它没有规定括号和运算符使用什么符号，也没有说明采用中缀、前缀还是后缀记法；没有说明如何区分三种基本表达式；甚至没有表示表达式第一项的 plus 可以省略，以及每项第一个因子的 $\times$ 必须省略。除去这种抽象程度和表示无关性，这个类型定义对应于一组 BNF 语法方程：

\quad $\langle$expression$\rangle ::= \langle$term$\rangle\ |\ \langle$addop$\rangle \langle$term$\rangle\ |\ \langle$expression$\rangle \langle$addop$\rangle \langle$term$\rangle$

\quad $\langle$term$\rangle ::= \langle$primary$\rangle\ |\ \langle$term$\rangle\langle$mulop$\rangle\langle$primary$\rangle$

\quad $\langle$primary$\rangle ::= \langle$unsigned real number$\rangle\ |\ \langle$variable$\rangle\ |\ (\langle$expression$\rangle)$

注意，我们尽可能用序列替代了递归。实际上，只要类型名在其定义的开头或末尾递归出现一次，就能这样做。例如：

\quad \textbf{type} expression $=$ \textbf{sequence} term;

\noindent
本可以递归地写成：

\quad \textbf{type} expression $= ($empty:$()$, non-empty:$($first:term; final:expression$))$.

\noindent
类似的替代表述，也能把 \textbf{while} 循环表达为递归过程。

构造递归定义类型的值，不需要新的运算符或转换函数；只需递归地使用为其他相关结构化方法定义的那些方法。例如，表达式

\quad $3/(b - 2)$

\noindent
可以用以下繁琐的构造指定：

\quad $[$term $($plus, $[$factor (times, primary $($const$(3)))$,

\tabto*{7.75em} factor $($div, primary $($bracketed $($

\tabto*{9em}  $[$term $($plus, $[$factor $($times, primary $($var$($``$b$''$)))]),$

\tabto*{9.1em}  term $($minus, $[$factor $($times, primary $($const$(2)))])])))$

\tabto*{7.55em}  $])$

\quad $]$.

\noindent
后面将在（\ref{sec:example-9-2}）中举例，介绍让计算机自己把表达式翻译为抽象结构的有效方法。

递归定义的数据的另一个熟悉例子，是家谱。家谱（不包括婚姻信息）可以这样定义：为每个人关联其所有子女各自的家谱。假定还需要保存一些个人详细信息：

\quad \textbf{type} family $= ($head:person; offspring:\textbf{sequence} family$)$;

没有子女的人，是家谱的最终分量，可以表示为：

\quad family$($Tom, $[])$

\noindent
有三个孩子的家庭可以表示为：

\noindent
family $($Jill, $[$family $($Tom, $[])$,

\tabto{3.2em} family $($Joanna, $[])$,

\tabto{3.2em} family $($Matthew, $[])])$.

最后一个例子说明，如何把 LISP 程序员熟悉的二叉树定义为递归数据结构。

\quad \textbf{type} list $= ($atom:\textbf{sequence} character, cons:$($car, cdr:list$))$.

\noindent
一个表在 LISP 的点记法中写成：

\quad $(($A.$($B.NIL$))$.NIL$)$

\noindent
它可以用几乎与 LISP 完全相同的方式，表达为 list 类型的值：

\quad cons $($cons $($atom $($``A''$)$,

\tabto{5em} cons $($atom $($``B''$)$, atom $($``NIL''$)))$,

\tabto{5em} atom $($``NIL''$)$

\tabto{4.8em} $)$;

\noindent
其中，转换为 list 类型的操作省略不写。

作为处理表的例子，我们编写一个函数，把整棵树翻转，使每个“左分支”变为“右分支”，反之亦然。

\quad \textbf{function} reverse$(l$:list$)$:list;

\quad \quad \textbf{with} $l$ \textbf{do}

\quad \quad \quad $\{$atom:reverse $\coloneq l$,

\quad \quad \quad ~ cons: reverse $\coloneq$ cons$($reverse$($cdr$)$, reverse $($car$))\}$

\subsection{表示}

递归类型的标准表示，与结构类似的非递归类型也很相近，区别在于：声明为递归类型自身的每个分量，不直接保存其值，而是用一个保存指向该值的指针的存储单元来表示。采用指针，是因为分量值的大小可能任意，无法分配固定数量的存储空间来容纳它。这称为“树形表示”，类似于数组或序列的树形表示，只是各分支可以增长到任意且不同的高度。

另一种表示方法是线性序列或\textit{位流}（bitstream）。这种表示可以不用指针，把递归子结构的值与其余信息连续存放，就像熟悉的、带括号的表达式字符表示一样。不过，我们可以根据上下文恢复括号结构，必要时扫描联合值的判别标记，而不用实际括号。这种方法通常与其他分量的紧缩表示结合使用，可以显著减少存储空间，代价是必须顺序访问结构中的各个分量。在许多情况下，位流表示的大小大约只有树形表示的十分之一。

\begin{figure}[h]
	\begin{subfigure}[t]{\textwidth}
		\centering
		\includegraphics[width=.7\textwidth]{chap2/fig8a}
		\caption{}
	\end{subfigure}
	
	\begin{subfigure}[t]{\textwidth}
		\centering
		\includegraphics[width=.58\textwidth]{chap2/fig8b}
		\caption{}
	\end{subfigure}
	\caption{(($A$\,.\,($B$.NIL))\,.\,($B$\,.\,$A$) 的表示}
\end{figure}

树形表示与线性表示之间的选择通常很明显。如果程序正在处理结构，通常是通过递归过程处理，那么方便访问任一分量的需要，就决定了应采用树形表示。此外，可以让多棵树共享公共分支，收回部分空间开销；共享分支是符号表达式处理的一个特点。但如果结构需要输出，随后重新输入，线性结构就优越得多。它不仅通过减少数据量缩短传送时间，线性化还避免了在辅助存储器上表示指针的许多棘手问题。许多情况下，一个结构经历多个处理和输入输出阶段，会在每个阶段转换表示；高级程序设计语言的多遍翻译程序就普遍这样做。

必须注意，共享递归子结构只是一种节省时间和存储空间的手段，对程序运行没有影响。这意味着，只要共享子结构有可能作为某个所属结构的一部分而被选择性更新，就必须避免共享。原则上，所有值彼此完全分离，程序员既无法知道，也无须关心自己的结构有多少部分被共享。此外，绝不可能让一个指针指回它所属的结构，因为这意味着该结构与自身的一个分量相同。只有无限结构才可能具有这种性质；而无限结构不满足排除公理，递归结构的重要归纳原理正是以这一公理为基础的。

\subsection{例子}
\label{sec:example-9-2}

程序设计语言中表达式的源文本，以如下定义的符号序列给出：

\quad \textbf{type} symbol $= ($constant:$($value:real$)$, variable:$($identifier: ident$)$,

\tabto*{9.2em}op:operator, leftbracket, rightbracket$)$;

编写一个程序，操作输入变量

\quad source:\textbf{sequence} symbol,

\noindent
从其开头读取尽可能长的合法表达式，返回相应的抽象表达式；如果无法做到，就转到标号 error 退出。结果的结构和源文本的语法，如本节前面所述。

程序结构紧密遵循所需结果的结构。有三个函数：

\quad compile expression

\quad compile term $($sign$)$

\quad compile primary

\noindent
每个函数都从 source 中移除其语法类别下最长的表达式，并返回相应的抽象结构。主要的特殊之处是，表达式的第一项可以没有符号；因此，sign 作为参数传给 compile term，而不是由 compile term 自己从 source 中读取。每个函数成功时，都有缩短 source 序列的副作用；失败时，则跳转到 error。

\noindent
\textbf{function} compile expression:expression;

\noindent
\quad \textbf{begin} sign:operator;

\noindent
\quad \quad \textbf{if} source.first $=$ plus $\vee$ source.first $=$ minus \textbf{then} sign \textbf{from} source

\noindent
\quad \quad \textbf{else} sign $\coloneq$ plus;

\noindent
\quad \quad compile expression $\coloneq [$compile term$($sign$)]$;

\noindent
\quad \quad \textbf{while} source.first $=$ plus $\vee$ source.first $=$ minus \textbf{do}

\noindent
\quad \quad \quad \textbf{begin} sign \textbf{from} source;

\noindent
\quad \quad \quad \quad compile expression$\wideparen{\text{:\ }[}$compile term$($sign$)]$

\noindent
\quad \quad \quad \textbf{end}

\noindent
\quad \textbf{end};

\noindent
\textbf{function} compile term$(s$:operator$)$:term;

\noindent
\quad \textbf{begin} $p$:primary; sign:operator; $fs$:\textbf{sequence} factor;

\noindent
\quad \quad $p\coloneq$ compile primary;

\noindent
\quad \quad $fs \coloneq [$factor(times, $p)]$;

\noindent
\quad \quad \textbf{while} source.first $=$ times $\vee$ source.first $=$ div \textbf{do}

\noindent
\quad \quad \quad \textbf{begin} sign \textbf{from} source:

\noindent
\quad \quad \quad \quad $p\coloneq$ compile primary;

\noindent
\quad \quad \quad \quad $fs\wideparen{\text{: }[}$factor$($sign, $p)]$

\noindent
\quad \quad \quad \textbf{end};

\noindent
\quad \quad compile term $\coloneq$ term$(s$, $fs)$

\noindent
\textbf{function} compile primary:primary;

\noindent
\quad \textbf{begin} $s$:symbol;

\noindent
\quad \quad $s$ \textbf{from} source;

\noindent
\quad \quad \textbf{with} $s$ \textbf{do} $\{$constant:compile primary $\coloneq$ const$($value$)$,

\noindent
\quad \quad \quad variable:compile primary $\coloneq$ var$($identifier$)$,

\noindent
\quad \quad \quad leftbracket:

\noindent
\quad \quad \quad \textbf{begin from} source;

\noindent
\quad \quad \quad \quad compile primary $\coloneq$ bracketed$($compile expression$)$;

\noindent
\quad \quad \quad \quad $s$ \textbf{from} source;

\noindent
\quad \quad \quad \quad \textbf{if} $s \neq$ rightbracket \textbf{then go to} error

\noindent
\quad \quad \quad \textbf{end},

\noindent
\quad \quad \textbf{else go to} error$\}$

\noindent
\quad \textbf{end};

\noindent
\textit{练习}
\nopagebreak

\noindent
编写程序，把表达式从树形表示转换为位流，再转换回来。
\section[稀疏数据结构]{稀疏数据结构}
\label{sec:sparse-data-structures}

讨论数组和幂集的表示时，迄今一直假定幂集的基类型和数组的定义域类型比较小，可以为结构中每个潜在元素分配一个位，或更大的存储区域来保存其值。例子也限于这类情况。本节研究基类型或定义域类型很大乃至无限、因而标准表示不再可行时，会带来什么后果和问题。

幂集及无限定义域映射，仍可以表示和操纵，条件是只考虑含有限个元素的集合，以及只有有限个元素取有意义值的映射；这里，“有意义”定义为不同于某个指定的空值或默认值。无限集合的幂集显然也是无限的；但由于这里幂集类型的每个值都只含有限个元素，只需在有限时间内列出这些元素，就能指定每个值，而列表也只占有限存储空间。类似地，对于定义域无限的映射类型，可以列出定义域中所有映射到有意义值的元素，并同时列出各自映射到的值域值，从而有限地指定其每个值。受到这种限制的类型，称为\textit{稀疏}类型。

实际上，稀疏性并不限于无限的基类型和定义域；它也适用于很大但有限的幂集，前提是程序员知道，自己关心的每个实际集合只会包含潜在元素中极小的一部分。例如，基类型可能包含数亿个值，但程序员知道只需处理大小不足一百的集合，而其中大多数甚至不足十。位模式表示会需要数亿个位，不可能采用；但由于程序中实际使用的每个值只有少数元素，用相对很少的空间就能方便地列出它们。这种幂集类型称为\textit{稀疏}幂集类型。类似地，定义域很大、其中几乎所有元素都映射到同一个默认值域值的数组，被称为属于\textit{稀疏数组}类型。

高级数据处理应用经常遇到稀疏集合和数组，其表示和操纵会带来一些熟悉的问题。第一个例子，是定义以车牌号码集合为值的类型。carnumber 类型的基数可能大约是四十亿；但程序员只想处理单个人所拥有汽车的集合，其中大多数只有一个元素，超过十个的极少。因此，carset 类型可以声明为稀疏幂集：

\quad \textbf{type} carset $=$ \textbf{sparse powerset} carnumber;

作为稀疏数组的例子，可以考虑从车主到其所拥有汽车集合的映射类型。每位车主以姓名和地址表示；两者长度任意，所以 owner 类型可以定义为：

\quad \textbf{type} owner $=$ \textbf{sequence} character;

\noindent
它的基数无限。因此，所需类型声明为稀疏类型：

\quad \textbf{type} carfile $=$ \textbf{sparse array} owner \textbf{of} carset.

在数据处理应用中，carfile 类型的变量称为\textit{随机访问}文件，而 owner 称为文件的\textit{键}元素。

接下来的两个例子来自数值应用。向量是从整数到浮点数的映射。稀疏向量的多数元素为零；因此，它的初值是常值零函数，所有元素都保持为零，除非显式赋予其他值：

\quad \textbf{type} sparsevector $=$ \textbf{sparse array} integer \textbf{of} real.

\noindent
稀疏复数矩阵可以用类似方式定义：

\quad \textbf{type} irregular matrix $=$ \textbf{sparse array} $($row, column:integer$)$ \textbf{of} complex.

下一个例子来自把程序设计语言翻译成机器码的工作。翻译过程中，翻译程序需要知道源程序中每个已声明标识符的某些信息，例如为变量分配的机器地址、长度、类型等等。假定这些信息属于 decode 类型。为每个标识符关联一个 decode 值的数组类型，命名为 dictionary，声明如下：

\quad \textbf{type} dictionary $=$ \textbf{sparse array} ident \textbf{of} decode

\noindent
当然，翻译程序只关心源程序中实际声明的那些标识符的 decode。对于绝大多数可能的标识符，这个类型的任何 dictionary 都会给出 decode 类型中表示“标识符未声明”的那个值。

最后一个例子，是商业文件系统和现实生活中都会引发熟悉问题的一种类型\textemdash{}多维交叉分类。一家公司的客户按地域分为若干区域，也按所购产品类别分为若干类别。有时需要访问某一区域的所有客户，依次遍历各类别；有时需要访问某一类别的所有客户，依次遍历各区域；最后，有时需要处理某一区域中某一类别的所有客户。处理这种情况所需的抽象结构，是一个以稀疏集合为元素的二维稀疏数组：

\quad \textbf{sparse array} $(c$:class; $a$:area$)$ \textbf{of sparse powerset} customer.

描述人与人的家庭关系时，也可能出现类似例子：

\quad \textbf{type} children $=$ \textbf{sparse array} $($mother, father:person$)$ \textbf{of} 

\tabto*{9.2em}\textbf{sparse powerset} person;

\noindent
对于涉及多次婚姻的情况，这个数组比可以定义为递归结构的、更接近树形的家庭表示更合适。

对于稀疏数组，有时把它们看作偏映射，而不是全映射，会很有用。偏映射不一定为定义域类型的每个元素都给出一个值。换言之，它实际有定义的定义域，是定义域类型的一个子集。对于这种数组类型，需要引入附加常量 omega，表示处处未定义的映射。再引入函数

\quad domain$(x)$

\noindent
也很有用，它返回一个集合，包含 $x$ 中实际有定义的元素的下标。因此，程序员可以遍历所有已定义的元素，或测试某个具体元素是否有定义。上述许多例子也完全可以声明为偏映射，而不是稀疏映射。对于偏映射，不必记录默认值。

\subsection{表示}

\noindent
稀疏集合和数组的通常表示，只记录默认值和那些有意义的元素。因此，这种表示借用序列类型处理大小可变结构的技术。稀疏集合可以看作稀疏映射的特例：把所有集合元素映射为布尔值 true，把所有非集合元素映射为默认值 false。因此，它们的表示与稀疏数组非常相似，不必单独讨论。

稀疏映射由若干元素组成。映射的每个元素都表示为其下标与值的笛卡尔积；这里，下标称为\textit{键}，值称为与该元素关联的\textit{信息}，两者并置在一起就称为一个\textit{条目}。对于稀疏集合，不需要记录任何信息，因为键本身的存在，足以表明该值是集合的元素。因此，稀疏集合的条目只包含键。

\subsubsection{顺序表示}

稀疏数组类型最简单的表示，是条目的序列；
即：

\quad \textbf{sparse array} $D$ \textbf{of} $R$

\noindent
表示得仿佛原先声明为：

\quad $($default:$R$; $s$:\textbf{sequence} $($key:$D$; information:$R))$.

现在，必须按选择序列表示时的同样准则，从各种可能的序列表示中选出一种。但用序列表示稀疏数组时，条目顺序无关紧要，无须反映它们加入的先后时间。因此，常常把条目按其
键值排序，尤其是当以后也将按这个顺序扫描时。

顺序表示的主要缺点，是访问任意给定下标对应的元素所需时间很长。对于规模较大的结构，程序设计者通常会花相当大的力气，确保按条目在序列中存储的同一标准顺序访问它们；如果要插入新条目，也先排序，再与原序列归并。因此，商业应用中标准的批处理和顺序文件更新方法，可以视为在不太适合随机访问的磁带介质上，实现稀疏数组这一抽象概念的实用办法。

\subsubsection{表式表示}

如果稀疏映射的条目个数有一个足够小、可以接受的上限 $N$，表式表示就能显著提高查找速度。其中，稀疏映射

\quad \textbf{sparse array} $D$ \textbf{of} $R$

\noindent
表示为非稀疏数组：

\quad $($default:$R$; occupied:\textbf{powerset} $0\,.\,.\,N$;

\tabto*{8em}\textbf{array} $0\,.\,.\,N$ \textbf{of} $($key:$D$; information:$R))$.

\noindent
如果所有有意义的条目都在使用之前收集完毕，就可以把表排序，然后用对数时间的查找，快速找到给定键对应的条目。

如果访问数组元素与添加新条目交替进行，就适合采用某种散列表技术。为此，选定一个任意的“散列”函数，把定义域类型 $D$ 映射为 $0\,.\,.\,N$ 范围内的整数。插入条目时，把它放在表中这个位置；以后访问它时，使用同一散列函数就能在这里找到它。如果该位置已经被另一个键不同的条目占用，就必须在表中寻找其他空位置。通常会搜索表中紧接着的后续位置；但表快满时，这可能使条目在某个恰好热门的区域周围不理想地聚集。解决办法之一，是选择 $N + 1$ 为素数，并用第二个散列函数从给定键算出任意步长。如果某位置已满，就把这个步长加到前一个位置上（模 $N + 1$），得到下一个要尝试的位置。

\subsubsection{索引表示}

表式存储方法只适合整张表都能
容纳在计算机主存中的情况。通常，如果
无法做到这一点，就把表式与顺序方法
结合使用。此时，把稀疏数组表示为一个表，其中每个
条目都是序列：

\quad $($default:$R$; table:\textbf{array} $1\,.\,.\,N$ \textbf{of}

\tabto*{10em}$($max:$D$; seq:\textbf{sequence} $($key:$D$; information:$R)))$.

\noindent
每个条目都放在这样一个序列 $i$ 中：其键落在 table $[i - 1]$.max（$i = 1$ 时为 $D$.min）与 table $[i]$.max 之间。表按序排列，以便快速定位合适的序列。这种技术类似多卷百科全书的编排：每卷书脊上注明该卷第一个和最后一个条目的键，不必把书从书架上取下，就能迅速识别正确的卷。

采用这种表示时，最好保证各序列长度大致相同。特别是，如果使用磁盘辅助存储器，让每个序列都放在单个柱面内非常有利，因为一次随机访问就不会需要多于一次磁头移动。因此，当某个序列过长时，它必须与相邻序列交换内容。这需要提取键最大和/或最小的条目；所有序列都已按键值排序时，最容易完成。排序和重新安排常常作为独立操作，定期进行；这种文件组织的总体方法称为“索引顺序”组织。

自然，在这种表示方法中，序列应尽可能短，例如短于磁盘上的一个磁道。因此，表本身可能变得很大，主存不再能容纳。这时，把表本身也分成若干段，再次使用同一原则，建立一个指向各段的第二级表。因此，每次访问数组元素，通常至少需要两次辅助存储访问。强烈建议，选择序列和表段的大小与位置时，应使它们与存储介质的访问特性紧密对应。

\subsubsection{局部稠密表示}

数值计算应用中遇到的一种特殊稀疏数组，是稀疏矩阵。稀疏矩阵常常可以划分为多个子矩阵，其中只有少数包含有意义的非零条目。这种情况下，可以说矩阵是局部稠密的，表示和处理时应利用这一事实。

一种方法，是为每个有意义的子矩阵保存其位置和大小，并把整个矩阵表示为这些子矩阵的表或序列；每个子矩阵按通常方式连续存储，使用乘法计算地址。不过，各子矩阵大小一般不同；如果处理矩阵时大小还会变化，问题就很严重。一种可能的方法，是把稀疏矩阵分成标准大小的子矩阵，例如十六乘十六，并建立一个表，保存指向各子矩阵的指针。全零子矩阵用空指针表示，不占其他存储空间；否则，按通常方式存储子矩阵，采用下面的地址计算方法。

每次访问数组时，先把两个下标的位值“交错”排列，使结果的最低有效部分包含两个下标各自的最低有效部分。随后，用结果的较高有效部分查地址表，定位所需子矩阵；用较低有效部分找到所需元素在子矩阵中的位置。某些机器上，这种下标交错技术可能比一般乘法更高效。如果某些子矩阵必须保存在辅助存储器中，就特别推荐这种地址计算方法，因为它按行处理矩阵与按列处理同样高效。因此，对于所有大型数组，无论是否稀疏，都可以推荐这种方法，尤其是在采用分页的计算机上。该方法的发明者是 E. W. Dijkstra 教授。

\subsubsection{网格表示}

文件交叉分类，在计算机上引起的问题与现实生活中一样多。通常的解决办法，是把最方便的一种分类作为标准，接受按另一种分类处理的额外成本，即使这需要重新排序文件。因此，稀疏映射

\quad \textbf{sparse array} $(i:D_1; j:D_2)$ \textbf{of} $R$

\noindent
表示为：

\quad \textbf{sparse array} $D_1$ \textbf{of} $($\textbf{sparse array} $D_2$ \textbf{of} $R)$

不过，也可以采用基于序列链式表示的方法，更对称地处理两个维度。在这种表示中，把 $D_1$ 的每个实际使用值放入一条链，把 $D_2$ 的每个实际使用值放入另一条链；这两条链称为\textit{边界链}。任一边界链的每个元素，都包含一个基址单元，指向具有该分类键值的所有元素组成的链式序列。现在，数组的每个实际条目都附有\textit{两个}地址：一个指向按 $D_1$ 分类时属于同一类的序列中的下一项，另一个指向按 $D_2$ 分类时属于同一类的下一项。因此，可以把每个项想象成位于二维网格线的一个交点，指针分别横向和纵向指向同一行或同一列的下一项。

\begin{figure}[h]
	\centering
	\includegraphics[width=.7\textwidth]{chap2/fig9}
	\caption{$A$: \textbf{sparse array} $(d_1: D_1;\ d_2: D_2)$ \textbf{of} $T$ 的网格表示}
\end{figure}

遗憾的是，这种网格表示只适合整个结构都能放入主存的情况。如果序列的大部分必须保存在辅助存储器中，为了效率，就应对相邻元素采用某种分块方法。
\section[例子：考试时间表]{例子：考试时间表}
\label{sec:example-examination}

在允许学生自由选择多种课程组合的教育机构中，需要编排这样一份考试时间表：每门考试只安排在一个时段，同时每位学生都能参加自己所修各门课程的考试。为每门考试单独分配一个时段，总能做到这一点；但无论考官还是学生，都希望整个考试期尽可能短。这意味着，每个时段应尽可能多地安排考试，但不能超过某个上限 $k$。考场大小还带来另一个约束：它只能容纳一定数量的学生。

在设计程序之前，最好根据各种已给定及所需数据项的结构，作更严格的形式化描述，以确认我们对问题的理解。“student”和“exam”类型显然是非结构化的，现阶段无须进一步定义。每位学生需要参加的考试，由一个映射给出：

\quad load:\textbf{array} student \textbf{of powerset} exam.

\noindent
时间表是时段的集合，而每个时段又由考试的集合组成：

\quad \textbf{type} session $=$ \textbf{powerset} exam;

\quad timetable:\textbf{powerset} session.

\noindent
接下来，尝试把输入和输出数据应具备的性质形式化。

\begin{enumerate}[wide, nosep, label=(\arabic*)]
	\item \label{en:examination-1}
	我们决定不把“时段个数最少”这一条件形式化，因为如果计算绝对最小值代价过高，我们实际上并不要求绝对最小值。

	\item \label{en:examination-2}
	每门考试都安排在某个时段：

	$\displaystyle\bigcup_{s\text{ \textbf{in} timetable}} s =$ exam.all
	\smallskip

	\item \label{en:examination-3}
	任何考试都不安排在多个时段：

	\quad $(s1 \wedge s2 = \{\})$ 或 $(s1 = s2)$

	\noindent
	条件 \ref{en:examination-2} 和 \ref{en:examination-3} 实际上说明，时间表把全部考试的集合划分为一组覆盖全部且互不相交的子集。

	\item \label{en:examination-4}
	任何时段都不包含超过 $k$ 门考试：

	\quad $s$ \textbf{in} timetable $\supset$ size$(s) \leqslant k$

	\item \label{en:examination-5}
	任何时段涉及的学生人数都不超过 hallsize。要形式化这一点，需要统计参加每门考试的学生人数：
	
	\quad examcount$(e$:$\text{exam}) = $ size $\{st$:student $\mid$ $e$ \textbf{in} load$(st)\}$.
	
	于是，一个时段涉及的学生人数是：
	
	\quad session count$(s\text{:session}) = \displaystyle\sum_{e \text{ \textbf{in} } s}$ examcount$(e)$
	
	条件可以形式化为：
	
	\quad $s$ \textbf{in} timetable $\supset$ sessioncount$(s) \leqslant$ hallsize.
	
	\item \label{en:examination-6}
	任何学生都不会在一个时段参加多门考试。要形式化这一点，引入考试不相容的概念：如果某个学生同时需要参加两门考试，这两门考试就不相容。对于每门考试 e1，都有一个与它不相容的考试集合 incompat$(e1)$：
	
	\quad incompat$(e1)= \{e2\text{:exam}\ |\ e2 \neq e1\ \&\ \exists\ st$:student$(e1 \text{ \textbf{in}}$ load$(st)$
		
	\tabto*{8.2em} \& $e2$ \textbf{in} load$(st))\}$
\end{enumerate}

\noindent
现在可以规定，一个时段中每一对考试都必须相容：

\quad $s$ \textbf{in} timetable \& $e1$, $e2$ \textbf{in} $s \supset\neg e1$ \textbf{in} incompat$(e2)$.

这六个条件用 load、hallsize 和 $k$ 定义。现实中任何成功的时间表，以及时间表的任何成功的计算机表示，都必须满足它们。它们规定了考试排程程序的目标与正确性标准。

\subsection{抽象程序}

考察这些条件可以发现，构造时间表不需要完整了解每位学生的 load。只需要每门考试的 examcount，以及与它不相容的其他考试的集合：

\quad examcount:\textbf{array} exam \textbf{of} integer;

\quad incompat:\textbf{array} exam \textbf{of powerset} exam.

\noindent
这两个数组体现了对实际数据的抽象，恰好把注意力集中在当前目的所需要的特征上，让我们暂时忽略情境中的其他特征。显然，只扫描一次学生的 load 数据，就很容易构造出这两个数组：

\quad examcount $\coloneq$ all$(0)$;

\quad incompat $\coloneq$ all$(\{\})$;

\quad \textbf{for} $st$:student \textbf{do}

\quad \quad \textbf{for} $e$ \textbf{in} load$(st)$ \textbf{do}

\quad \quad \quad \textbf{begin} examcount$(e)$: $+\ 1$;

\quad \quad \quad \quad incompat$(e)$: $\vee\ (\text{load}(st) - \{e\})$

\quad \quad \quad \textbf{end};

寻求解法时，一个简化因素是这些条件很容易分成两类：\ref{en:examination-1}、\ref{en:examination-2} 和 \ref{en:examination-3} 针对整个时间表，而 \ref{en:examination-4}、\ref{en:examination-5} 和 \ref{en:examination-6} 只针对单个时段，根本不涉及整个时间表。因此，程序可以组织为一个内层部分，选出满足 \ref{en:examination-4}、\ref{en:examination-5} 和 \ref{en:examination-6} 的合适时段；再加一个外层循环，用这些合适时段构造时间表。

外层循环的目标，是在完成时满足条件 \ref{en:examination-2} 和 \ref{en:examination-3}。因此，选择其中一个作为循环终止条件，并把循环体设计为始终保持另一个条件成立（即循环不变式）；此外，还要保证开始循环之前，不变式已经成立。

显然可以选择互不相交性（条件 \ref{en:examination-3}）作为不变式，把覆盖全部作为终止条件，使每次执行循环体都向它推进。空时间表显然满足不变式。这样就得到如下结构的算法：

\quad timetable $\coloneq \{\}$;

\quad \textbf{while} timetable does not satisfy \ref{en:examination-2} \textbf{do}

\quad \textbf{begin} select a session satisfying \ref{en:examination-4}, \ref{en:examination-5}, \ref{en:examination-6};

\quad \quad add the session to the timetable

\quad \textbf{end};

\quad print timetable.

要使添加新时段后仍保持
不变式成立，该时段的考试必须从
尚未出现在时间表中的考试里选取。因此，引入一个
新变量保存这些剩余考试：

\quad remaining:\textbf{powerset} exam;

\noindent
它由如下不变关系定义：

\quad remaining $=$ exam.all $- \displaystyle\bigcup_{s \text{ \textbf{in} timetable}} s$.

\noindent
现在，整个程序的结构成为：

\quad timetable $\coloneq \{\}$;

\quad remaining $\coloneq$ exam.all;

\quad \textbf{while} remaining $\neq \{\}$ \textbf{do}

\quad \quad \textbf{begin} $s\coloneq$ suitable;

\quad \quad \quad timetable: $\vee\ \{s\}$;

\quad \quad \quad remaining: $-\ s$

\quad \quad \textbf{end};

\quad print timetable.

剩下的问题，是每个阶段如何选出合适时段。原则上，没有理由认为每一步的“最佳”选择，最终会产生“最佳”甚至“良好”的时间表。不过，一般看来，选择剩余考试的一种组合，使它尽量接近填满考场，或尽量接近上限 $k$，应该有利。这样，可能用一组相当紧凑的时段安排大多数学生和考试，尽管最后可能留下一个颇长的“尾巴”，包含少数学生参加的许多小型时段。虽然这不会使时段个数最少，但对大多数学生和考官可能已经相当满意。

如果不接受每一步表面上的最佳选择，就只能尝试更全局的优化；这要么涉及天文数字般的试探次数，要么需要一些复杂考虑，而这些考虑很可能只有在积累了较简单算法的实际经验之后才会显现。因此，现阶段除了希望取得好结果，也没有别的办法。

还需要编写函数：

\quad \textbf{function} suitable:session,

\noindent
从剩余考试集合中选出合适时段。一种可能的方法，是生成若干满足 \ref{en:examination-4}、\ref{en:examination-5} 和 \ref{en:examination-6} 的试探时段，选出找到的最佳者。最佳时段很可能是 sessioncount 最大的那个；但以后可能要调整选择标准，所以最好把选择过程定义为独立子程序，用于更新变量

\quad bestsofar:session;

\noindent
更新依据是以下变量的当前值：

\quad trial:session;

\quad \textbf{procedure} record

\quad \quad \textbf{if} sessioncount$($bestsofar$) <$ sessioncount $($trial$)$ \textbf{then}

\quad \quad \quad bestsofar $\coloneq$ trial.

\noindent
suitable 的结果就是 bestsofar 的最终值：

\quad suitable $\coloneq$ bestsofar.

还需编写一个过程，生成并记录一系列满足 \ref{en:examination-4}、\ref{en:examination-5} 和 \ref{en:examination-6} 的试探时段。考察这些条件可知，如果某个试探时段\textit{不}满足其中一个条件，更大的试探时段也不会满足它。换言之，找到一个\textit{不可行}的试探时段以后，不必再生成包含它的其他试探时段。因此，可以把生成过程组织为：生成已知可行的每个试探时段的所有超集，但排除已经尝试过的考试。为此，引入变量：

\quad untried:\textbf{powerset} exam,

\noindent
以及过程：

\quad \textbf{procedure} gensupersets,

\noindent
它从 untried 中给 trial 添加一个或多个考试，生成并记录 trial 的所有可行超集。这个过程将在“suitable”内部调用。

\quad \textbf{function} suitable:session;

\quad \quad \textbf{begin} trial, bestsofar:session; $e$:exam; untried:\textbf{powerset} exam;

\quad \quad \quad $e$ \textbf{from} remaining;

\quad \quad \quad trial $\coloneq$ bestsofar $\coloneq$ $\{e\}$;

\quad \quad \quad untried $\coloneq$ remaining $-$ trial $-$ incompat$(e)$;

\quad \quad \quad gensupersets;

\quad \quad \quad suitable $\coloneq$ bestsofar

\quad \quad \textbf{end};

\noindent
注意，trial 的初值是某门考试的单元素集合；该考试从剩余集合中选出，选择标准尚未确定。这样做的理由是，所选考试最终总要出现在时间表的某个时段中，不妨就放在当前这个时段。如果不先作这个选择，考试排程程序的主循环每轮都会让 gensupersets 不断生成相同的超集。

另一个重要优化，是从 untried 中移除所有与 trial 中考试不相容的考试，因为完全没有必要考虑把它们加入 trial。

生成给定 trial 的超集时，可以依次从 untried 中选出考试，加入 trial。如果结果仍有效，就记录它；trial 的新值此时就是一个合适时段，可以作为进一步生成超集的基础。这说明适合采用递归程序结构。当然，加入 trial 的考试还应从 untried 中移除，以免不必要的重复；同时，从 untried 中移除所有与刚加入考试不相容的考试，非常有利，因为以后不必再考虑它们。此外，每次调用结束后，trial 和 untried 的值必须保持原状；因此，对它们的任何修改，都必须借助 save 1 和 save 2 两个变量记录并恢复。

\textbf{procedure} gensupersets;

\quad \textbf{begin} $e$:exam; save 1, save 2:\textbf{powerset} exam;

\quad \quad record; save 1 $\coloneq$ untried;

\quad \quad \textbf{if} size$($trial$) < k$ \textbf{then}

\quad \quad \textbf{while} untried $\neq \{\}$ \textbf{do}

\quad \quad \quad \textbf{begin} $e$ \textbf{from} untried;

\quad \quad \quad \quad save $2\coloneq$ untried $\wedge$ incompat$(e)$;

\quad \quad \quad \quad untried:$-$ save 2;

\quad \quad \quad \quad trial:$\vee\ \{e\}$;

\quad \quad \quad \quad \textbf{if} sessioncount(trial) $<$ hallsize \textbf{then} gensupersets;

\quad \quad \quad \quad untried:$\vee$ save 2;

\quad \quad \quad \quad trial:$-\ \{e\}$

\quad \quad \quad \textbf{end};

\quad \quad \quad untried $\coloneq$ save 1

\quad \textbf{end} gensupersets.

这个程序的有效性取决于：trial 始终满足时间表中时段所需的条件 \ref{en:examination-4}、\ref{en:examination-5} 和 \ref{en:examination-6}，并且始终是 remaining 的子集。

理由如下：

对于 \ref{en:examination-4}：gensupersets 只有在 trial 的大小严格小于 $k$ 时，才生成超集。

对于 \ref{en:examination-5}：当 trial 的 sessioncount 大于考场容量时，绝不会进入 gensupersets（假定任何 examcount 都不大于 hallsize）。

对于 \ref{en:examination-6}：从 untried 中移除不相容集合，保证任何时刻 untried 中剩下的所有考试，都与 trial 中的所有考试相容。因此，把 untried 中任意一门考试转移到 trial，绝不会违反 \ref{en:examination-6}。

最后，在初次调用 gensupersets 时，untried $\subset$ remaining。untried 本质上是一个不增加的量：每次向它添加元素，之前都已经移除了同一批元素。因此，untried 始终是 remaining 的子集；而 trial 只由 untried 的元素构造，所以它也必须始终是 remaining 的子集。

这样，就完成了构造考试时间表的抽象程序的第一个版本。把所有部分汇集起来，如下所示：

hallsize, $k$:integer, initially given;

load:\textbf{array} student \textbf{of powerset} exam, initially given;

\textbf{type} session $=$ \textbf{powerset} exam;

timetable:\textbf{powerset} session, initially $\{\}$;

examcount:\textbf{array} exam \textbf{of} integer, initially all $(0)$;

incompat:\textbf{array} exam \textbf{of powerset} exam, initially constant $(\{\})$;

\textbf{function} sessioncount($s$:session):integer;

\quad \textbf{begin} sum:integer, initially 0;

\quad \quad \textbf{for} $e$ \textbf{in} $s$ \textbf{do} sum:$+$ examcount$(e)$;

\quad \quad sessioncount $\coloneq$ sum

\quad \textbf{end};

\quad remaining:\textbf{powerset} exam, initially exam.all;

\textbf{function} suitable:session;

\quad \textbf{begin} bestsofar, trial:session; untried:\textbf{powerset} exam;

\quad \quad $e$:exam; $e$ \textbf{from} remaining; bestsofar $\coloneq \{e\}$;

\quad \quad trial $\coloneq \{e\}$; untried $\coloneq$ remaining $-$ trial $-$ incompat$(e)$;

\quad \quad gensupersets;

\quad \quad suitable $\coloneq$ bestsofar

\quad \textbf{end};

\noindent
以下两个过程是 suitable 的局部过程：

\quad \textbf{procedure} record;

\quad \quad \textbf{if} sessioncount (bestsofar) $<$ sessioncount$($trial$)$ \textbf{then} 

\quad \quad \quad bestsofar $\coloneq$ trial;

\quad \textbf{procedure} gensupersets;

\quad \quad \textbf{begin} $e$:exam; save 1, save 2:\textbf{powerset} exam;

\quad \quad \quad record; save 1 $\coloneq$ untried;

\quad \quad \quad \textbf{if} size$($trial$) < k$ \textbf{then}

\quad \quad \quad \quad \textbf{while} untried $\neq \{\}$ \textbf{do}

\quad \quad \quad \quad \quad \textbf{begin} $e$ \textbf{from} untried;

\quad \quad \quad \quad \quad \quad save 2 $\coloneq$ untried $\wedge$ incompat$(e)$;

\quad \quad \quad \quad \quad \quad untried: $-$ save 2;

\quad \quad \quad \quad \quad \quad trial: $\vee\ \{e\}$;

\quad \quad \quad \quad \quad \quad \textbf{if} sessioncount$($trial$) <$ hallsize \textbf{then} gensupersets;

\quad \quad \quad \quad \quad \quad untried: $\vee$ save 2;

\quad \quad \quad \quad \quad \quad trial:$-\ \{e\}$

\quad \quad \quad \quad \quad \textbf{end};

\quad \quad \quad \quad untried $\coloneq$ save 1

\quad \textbf{end} gensupersets;

\noindent
主程序如下：

\quad \textbf{for} $st$:student \textbf{do}

\quad \quad \textbf{for} $e$ \textbf{in} load$(st)$ \textbf{do}

\quad \quad \quad \textbf{begin} examcount$(e)$: $+\ 1$; incompat$(e)$: $\vee\ ($load$(st) - \{e\})$ \textbf{end};

\quad \textbf{while} remaining $\neq \{\}$ \textbf{do}

\quad \quad \quad \textbf{begin} $s$:session;

\quad \quad \quad \quad $s\coloneq$ suitable;

\quad \quad \quad \quad timetable: $\vee\ \{s\}$;

\quad \quad \quad \quad remaining: $-\ s$

\quad \quad \quad \textbf{end};

\quad print timetable

在投入更多精力继续开发之前，最好严格审查这个程序，确保它能成功。现在，程序可能失败的最明显原因有：

\begin{enumerate}[wide, nosep, label=(\arabic*)]
	\item 最终时间表规模大到无法接受；我们已同意，在更加了解数据以前，对此无能为力。

	\item 每一步生成全部试探时段的耗时过多。在程序开始、remaining 仍很大时，这尤其严重；如果每次递归调用 gensupersets 时 untried 都仍很大，也会如此。减小 untried 的主要方法，是移除所有与 trial 不相容的考试。因此，应优先把不相容集合最大的考试加入 trial，使 untried 尽快缩小。如果按这个标准相同，就先选 examcount 最大的考试。这两个标准之间的具体权重，可能要根据以后的经验调整；目前，最简单的实现办法，是预先按这一标准排序考试，并让 $e$ \textbf{from} untried 选取第一个元素。
\end{enumerate}

如果这种简单策略仍然不够，可能必须人为限制 gensupersets 中循环的迭代次数。但要选择合适策略，恐怕需要实际经验；暂且希望不必如此。

\subsection{数据表示}

要设计成功的数据表示，需要知道
问题大概有多大。在这个例子中，我们作
以下假定：

\begin{enumerate}[wide, nosep, label=(\arabic*)]
	\item 考试不超过 500 门，每门参加人数少于 1000（典型值为 50）。
	\item 学生大约有 5000 人。
	\item 每位学生参加的考试少于十门，通常为五门。
	\item 考场大约容纳 1000 名学生。
	\item 可以接受的同时考试门数上限为 30，典型门数为 10。
	\item 人工排程方法曾成功编排出不超过 50 个时段的时间表。
\end{enumerate}

下面逐项考虑数据。

\begin{enumerate}[wide, nosep, label=(\arabic*)]
	\item \textbf{type} \textit{exam}

	\noindent
	显而易见的表示，是整数子界：$0\,.\,.\,500$。

	\item \textbf{type} \textit{session}

	\noindent
	显然有两种选择：位模式表示（500 位），或一个由 30 个九位元素组成的数组（$+\ \text{指针}$）（$270 \text{ 位 } + \text{ 一个字}$）。需要存储的时段个数不多，因此存储节省并不重要。对时段的主要操作，是插入一门已知尚未在其中的考试，以及移除最近插入的考试。因此，数组方法最好，因为可以用栈方法插入和移除元素。

	由于经常需要知道 sessioncount，把它与时段一起记录，并在插入、移除元素时同步维护，会很有利。

	trial 和 bestsofar 都采用这种表示。

	\item \textit{timetable}

	\noindent
	时间表上的唯一操作，是插入新时段。由于时段长度可变，可以把时间表组织为变长序列的序列。每门考试在时间表中恰好出现一次，因此时间表最大大小是 $500 \times \text{nine}$ 位，再加大约六十个字来指出各时段的分界（如果超过六十个时段，程序无论如何也算失败了）。

	另一种简单得多的表示，只记录每门考试属于哪个时段。它只需要：

	\quad \textbf{array} exam \textbf{of} $1\,.\,.\,60$

	\noindent
	这种表示之所以可行，完全是因为时间表各时段互不相交。

	\item \textit{examcount}:\textbf{array} \textit{exam} \textbf{of} \textit{integer}

	\noindent
	显然应选择标准表示。

	\item \textit{remaining}, \textit{untried}, \textit{save} 1, \textit{save} 2

	\noindent
	这些变量开始时相当满，随程序推进逐渐变空。因此，它们的平均密度约为百分之五十，采用稀疏表示没有意义。此外，对它们频繁执行标准集合操作，也说明应采用标准位模式表示。

	\item \textit{incompat}

	\noindent
	incompat 的元素最常见的用途，是从 untried 中减去它们，因此也应采用位模式表示。这将需要 $500 \times 500 \text{ 位}$，约为 10000 个字。它是所需数据结构中最大的一个，远大于其他结构；但它给程序带来的额外速度，可能足以证明其总体大小合理，而且多数要运行这个程序的计算机都能接受，所以似乎不值得寻求更紧凑的表示。

	\item \textit{load}

	\noindent
	每位学生的 load 是问题的主要输入数据，数据量也可能极大。因此，程序只需扫描一次数据，真是双重幸运：不仅可以把数据作为外部序列提供，还意味着其表示可以设计得适合人工阅读、书写和穿孔。

\end{enumerate}

因此，每位学生分配一张卡片，用十个各含六个字符的栏位保存考试编号。为减少不必要的穿孔，第一个空白栏表示考试集合结束。为了识别，每张卡还应包含学号；幸运的是，程序可以完全忽略它，不过可能还是应检查学号，以防重复或遗漏。

\noindent
\textit{练习}
\nopagebreak

\noindent
采用建议的数据表示，对上述抽象程序编码。
\section{公理化}

前面各节介绍了若干构造数据空间（类型）的方法，并说明了这些空间上定义的一些有用操作。但描述本质上仍是直观、非形式化的。因此，人们会问：有关数据空间的全部信息是否都已经说明，还是仍有可能误解某些细节？

为消除这种误解，或检查是否已经发生误解，最好为每个数据空间及其上定义的运算符，给出严格的数学规约。我们采用如今惯常的数学做法，不用传统定义，而用一组公理来严格规定推理所涉及的对象。

考虑到公理在数据结构化理论中的作用，概括其预期性质可能很有帮助。

\begin{enumerate}[wide, nosep, label=(\arabic*)]
	\item 公理形式化地陈述现实世界及其计算机内部表示共有的性质。正是凭借这些性质，计算机程序对表示的操纵才会产生可以成功应用回现实世界的结果。
	\item 公理建立一个概念框架，涵盖那些被认为与程序员任务有关的现实世界方面，从而帮助其构造和发明性的思考。
	\item 公理严格陈述计算机程序所依据的关于现实世界的假定。
	\item 公理陈述任何计算机数据表示都必须具备的性质，而不涉及初始阶段无关的细节。
	\item 公理给予程序员或高级语言实现者经过严格限定的自由，使其能选择最适合自己的应用和现有硬件的表示。
	\item 公理构成程序正确性的任何证明的基础。
\end{enumerate}

这里给出的公理，不打算直接用于证明复杂程序，因为这样的证明会过于冗长。它们更适合用来建立所描述数据空间的一些熟悉性质，然后在证明中非形式化地使用这些性质。最终，或许可以让计算机检查这类证明；但这需要发展比逻辑学家目前提供的语言强大得多的证明表达形式语言，并为定理的许多大型子类采用强大的判定过程，以帮助验证证明的各个步骤。

适用于给定类型的公理，取决于该类型的定义方式。因此，无法像整数公理那样，在每种情况下给出固定的一组公理；我们给出的是一种模式或图式，说明如何从相应类型定义的一般形式，导出一组具体公理。

\subsection{枚举类型与子界}

以下公理为枚举类型和子界所共有，仿照熟悉的自然数公理建立。假定类型名为 $T$，所有变量都属于这一类型。

\begin{enumerate}[wide, nosep, label=(\arabic*)]
	\item \label{en:enumerate-1}
	$T$.min 是 $T$ 类型的值
	
	\item \label{en:enumerate-2}
	若 $x$ 是 $T$ 类型的值，且 $x\neq T$.max

	\qquad~ 则 succ$(x)$ 是 $T$ 类型的值

	\item $T$ 的元素仅为 \ref{en:enumerate-1} 和 \ref{en:enumerate-2} 所规定的那些值

	\item $\text{succ}(x) = \text{succ}(y) \supset x = y$

	\item succ$(x) \neq T$.min

	\item $\text{pred}(\text{succ}(x)) = x$

	\noindent
	以下三条公理只适用于有序类型

	\setcounter{enumi}{6}
	\item $T$.min $\leqslant x$

	\item $x \leqslant T$.min $\supset x = T$.min

	\item \label{en:enumerate-9}
	$\text{succ}(x) \leqslant \text{succ}(y) \equiv x \leqslant y$

	\noindent
	注意：succ$(T$.max$)$ 和 pred$(T$.min$)$ 未定义。

	用枚举定义类型的一般形式是：

	\quad \textbf{type} $T = (k_1, k_2, \dots, k_n)$;

	\noindent
	其中 $T$ 是类型名，$k_1, k_2 , \dots, k_n$ 是该类型所有值的名称。
	
	该类型的附加公理是：
	
	\tabto*{.9em}(10) $k_1 = T$.min
	
	\tabto*{1.8em}\& $k_2 = \text{succ}(k_1)$
	
	\tabto*{1.8em}\& $k_3 = \text{succ}(k_2)$
	
	\tabto*{3.8em}\dots
	
	\tabto*{1.8em}\& $k_n = \text{succ}(k_{n-1}) = T$.max.
	
	把类型定义为子界的一般形式是：
	
	\quad \textbf{type} $T=k\,.\,.\,l;$
	
	\noindent
	其中 $k$ 和 $l$ 属于基类型 $T_0$。
	
	该类型的附加公理是：
	
	\tabto*{.9em}(10) $T$.min $= k$
	
	\tabto*{1.8em}\& $T$.max $= l$.

	\setcounter{enumi}{10}
	\item $T(T_0(x)) = x$.

	\item $k \leqslant x_0\ \&\ x_0 \leqslant l \supset T_0(T(x_0)) = x_0$.

	\item $x \leqslant y \equiv T_0(x) \leqslant T_0(y).$
\end{enumerate}

利用公理 \ref{en:enumerate-1} 到 \ref{en:enumerate-9}，可以证明以下次序性质：

\begin{enumerate}[wide, nosep, label=(T\arabic*)]
	\item $x \leqslant x$.

	\item \label{en:enumeration-axiom-2}
	$x \leqslant \text{succ}(y) \supset x = \text{succ}(y) \vee x \leqslant y$.
	
	\item \label{en:enumeration-axiom-3}
	$z \leqslant y\ \&\ y \leqslant x \supset z \leqslant x.$
	
	\item $x \leqslant y\ \&\ y \leqslant x \supset x = y.$
\end{enumerate}

\noindent
\textit{提示}：对 $x$ 归纳。证明 \ref{en:enumeration-axiom-2} 需要用到 \ref{en:enumeration-axiom-3}。

\noindent
简写：

\noindent
若 $\ominus$ 是一元运算符，$\oplus$ 是二元运算符，且两者都接受基类型 $T_0$ 的操作数，那么在 $a$ 属于 $T_0$ 类型、$x$ 和 $y$ 属于 $T$ 类型时，以下简写允许省略转换函数：

\begin{enumerate}[wide, nosep, label=(\arabic*)]
	\setcounter{enumi}{13}
	
	\item ~ $\ominus\ x$ 代表 $\ominus\ T_0(x)$。

	\item $x \oplus y$ 代表 $T_0(x) \oplus T_0(y)$。
	
	\item $x \oplus a$ 代表 $T_0(x) \oplus a.$

	\item $a \oplus x$ 代表 $a \oplus T_0(x).$

	\item $a \coloneq x$ 代表 $a\coloneq T_0(x)$。
\end{enumerate}

\subsection{笛卡尔积}

把类型定义为笛卡尔积的一般形式是：

\quad \textbf{type} $T = (s_1 : T_1;\ s_2: T_2; \dots ;\ s_n: T_n);$

\noindent
其中 $s_1, s_2, \dots , s_n$ 是分量的选择函数，$T_1, T_2, \dots, T_n$ 是相应分量的类型。

\begin{enumerate}[leftmargin=2\parindent, label=(\arabic*)]
	\item \label{en:cartesian-1}
	若 $x_1$ 是 $T_1$ 类型的值，$x_2$ 是 $T_2$ 类型的值，\dots，$x_n$ 是 $T_n$ 类型的值

	\quad 则 $T(x_1, x_2, \dots, x_n)$ 是 $T$ 类型的值。

	\item $T$ 的元素仅为 \ref{en:cartesian-1} 所规定的那些值。

	\item 若 $x = T(x_1, x_2, \dots, x_n)$，则

	\quad $x.s_1 = x_1\ \&\ x.s_2 = x_2\ \& \dots \&\ x.s_n = x_n$。

\end{enumerate}

\noindent
简写：

\begin{enumerate}[leftmargin=2\parindent, label=(\arabic*)]
	\setcounter{enumi}{3}
	\item $x.s_1 \coloneq x_1$ 代表 $x\coloneq T(x_1, x.s_2, \dots, x.s_n)$

	$x.s_2 \coloneq x_2$ 代表 $x\coloneq T(x.s_1, x_2, \dots, x.s_n)$

	\qquad \dots

	$x.s_n \coloneq x_n$ 代表 $x\coloneq T(x.s_1, x.s_2, \dots, x_n)$

	\item 若 $x$ 是 $T$ 类型的值，则
	
	\textbf{with} $x$ \textbf{do} $S$ 或 \textbf{with} $x$ \textbf{take} $S$ 代表
	
	\quad $S_{x.s_1, x.s_2, \dots, x.s_n}^{s_1, s_2, \dots, s_n}$
\end{enumerate}

\noindent
它表示：用 $S$ 的每个下标，替换 $S$ 中对应上标的全部自由出现。

\begin{enumerate}[leftmargin=2\parindent, label=(\arabic*)]
	\setcounter{enumi}{5}
	\item 在需要 $T$ 类型表达式的上下文中，$(x_1, x_2, \dots, x_n)$ 代表 $T(x_1, x_2, \dots, x_n)$。
\end{enumerate}

如果把笛卡尔积类型视为有序的，就适用以下公理：
	
\begin{enumerate}[leftmargin=2\parindent, label=(\arabic*)]
	\setcounter{enumi}{6}
	\item $x \leqslant y \equiv$\quad $x.s_1 < y.s_1$
		
	\tabto*{3.2em}$\vee\ x.s_1 = y.s_1\ \&\ (x.s_2 < y.s_2$
	
	\tabto*{3.2em}$\vee\ x.s_2 = y.s_2\ \&\ (x.s_3 < y.s_3$
	
	\tabto*{3.2em}$\vee\ \dots$ \tabto*{8.6em}$\&\ (x.s_{n - 1} < y.s_{n-1}$
	
	\tabto*{3.2em}$\vee\ x.s_{n-1} = y.s_{n-1}\ \&\ x.s_n \leqslant y.s_n)\dots)).$
\end{enumerate}

\subsection{带判别标记的联合}

定义的一般形式是：

\quad \textbf{type} $T = (s_1: T_1; s_2: T_2; \dots ; s_n:T_n; k_1: T'_1, k_2:T'_2, \dots , k_m: T'_m)$

\begin{enumerate}[wide, nosep, label=(\arabic*)]
	\item \label{en:distriminated-1}
	若 $x_1$ 是 $T_1$ 类型的值，$x_2$ 是 $T_2$ 类型的值，\dots，$x_n$ 是 $T_n$ 类型的值
	
	且 $x'_1$ 是 $T'_1$ 类型的值，$x_2$ 是 $T'_2$ 类型的值，\dots，$x'_m$ 是 $T'_m$ 类型的值
	
	则下列各值是 $T$ 中彼此不同的元素：
	
	$T(x_1, x_2, \dots, x_n, k_1(x'_1))$
	
	$T(x_1, x_2, \dots, x_n, k_2(x'_2))$
	
	\quad \dots
	
	$T(x_1, x_2, \dots, x_n, k_n(x'_n))$

	\item $T$ 的元素仅为 \ref{en:distriminated-1} 所规定的那些值
	
	\item\label{en:distriminated-3}
	对于 1 到 $m$ 之间的每个 $i$，若 $x = T(x_1, x_2, \dots, x_n, k_i(x'_i))$，则
	
	\qquad~ $x.s_1 = x_1\ \&\ x.s_2 = x_2\ \& \dots \&\ x.s_n = x_n$
	
	\quad~ \& $x.k_i = x'_i$

	\noindent
	\textit{注意}：\quad~~\, 当 $l \neq i.$ 时，$x.k_l$ 未定义。
	
	\noindent
	简写：
	
	\item 在与 \ref{en:distriminated-3} 相同的条件下，
	
	\quad~\, \textbf{with} $x$ \textbf{do} $\{k_1:S_1, k_2:S_2, \dots, k_n:S_n\}$ 表示

	\quad~\, \textbf{with} $x'_i$ \textbf{do} $(S_i)_{x_1, x_2, \dots, x_n}^{s_1, s_2, \dots, s_n}$;

	\quad~\, 用 \textbf{take} 替代 \textbf{do} 时同理。

	\item 若 $n = 0$，则 $k_i(x'_i)$ 代表 $T(k_i(x'_i)).$
\end{enumerate}

\subsection{数组}

数组定义的一般形式是：

\quad \textbf{type} $T =$ \textbf{array} $D$ \textbf{of} $R$

\begin{enumerate}[wide, nosep, label=(\arabic*)]
	\item \label{en:arrays-1}
	\textbf{若} $r$ 是 $R$ 类型的值，\textbf{则} $T(r)$ 是 $T$ 类型的值

	\item \label{en:arrays-2}
	\textbf{若} $x$ 是 $T$ 类型的值，$d$ 是 $D$ 类型的值，\textbf{且} $r$ 是 $R$ 类型的值

	\qquad\quad \textbf{则} $T(x, d:r)$ 是 $T$ 类型的值

	\item $T$ 的元素仅为 \ref{en:arrays-1} 和 \ref{en:arrays-2} 所规定的那些值。

	\item $T(T(x, d:r), d':r') =$

	\quad \textbf{if} $d = d'$ \textbf{then} $T(x, d':r')$

	\qquad\quad\textbf{else} $T(T(x, d':r'), d:r).$

	\item $T(r)[d] = r.$

	\item $T(x, d:r)[d'] =$ \textbf{if} $d' = d$ \textbf{then} $r$ \textbf{else} $x[d'].$

	\item $($\textbf{for} $i:D$ \textbf{take} $E(i))[j] = E(j).$

	\noindent
	简写：

	\item $x[d] \coloneq r$ 表示 $x \coloneq T(x, d:r)$。

	\item\label{en:arrays-9}
	 $T(x, d_1:r_1, d_2 : r_2, \dots, d_n:v_n)$ 代表

	\quad $T(T(\dots T(T(x), d_1:v_1), d_2:v_2)\dots), d_n:r_n)$

	\item 在 \ref{en:arrays-9} 中，如果 $d_1, d_2, \dots, d_n$ 遍及整个定义域类型，就可以省略 $x$。类似地，在合适的上下文中，也可以省略 T。

如果数组类型有序，就适用以下公理：

	\item $x \leqslant y \equiv \forall d: D(y[d] < x[d] \supset \exists d':D(d' < d\ \&\ x[d'] < y[d']))$
\end{enumerate}

\noindent
定理：

$x = y \equiv \forall d:D(x[d] = y[d])$

\subsection{幂集}

以下集合公理只适用于\textit{有限}集合，而且这些集合的类型必须按层级组织。因此，可以避开那些威胁更强集合论版本的公理化的悖论。

幂集定义的一般形式是：

\quad \textbf{type} $T =$ \textbf{powerset} $T_0$,

\noindent
其中 $T_0$ 是基类型。

\noindent
设 $a$、$b$ 是 $T_0$ 类型的值。

\begin{enumerate}[wide, nosep, label=(\arabic*)]
	\item \label{en:powerset-1}
	$T()$ 是 $T$ 类型的值
	\item \label{en:powerset-2}
	\textbf{若} $x$ 是 $T$ 类型的值，\textbf{且} $a$ 是 $T_0$ 类型的值，\textbf{则} $x \vee T(a)$ 是 $T$ 类型的值
	\item $T$ 的元素仅为 \ref{en:powerset-1} 和 \ref{en:powerset-2} 所规定的那些值。
	\item $\neg a$ \textbf{in} $T()$
	\item $a$ \textbf{in} $(y \lor T(a))$
	\item $a \neq b \supset (a \text{ \textbf{in} } (x \lor T(b)) \equiv (a \text{ \textbf{in} } x)$
	\item $T() \subset x$
	\item $(y \lor T(a)) \subset x \equiv (y \subset x\ \&\ a \text{ \textbf{in} } x)$
	\item $x = y \equiv (x \subset y)\ \&\ (y \subset x)$
	\setcounter{enumi}{9}
	\item $x \lor T() = x$
	\item $x \lor (y \lor T(a)) = (x \lor T(a)) \lor y$
	\item $x \land T() = T()$
	\item $x \land T(a) =$ \textbf{if} $a$ \textbf{in} $x$ \textbf{then} $T(a)$ \textbf{else} $T()$
	\item $x \land (y \lor T(a)) = (x \land y) \lor (x \land T(a))$
	\item $T() - x = T()$
	\item $T(a) - x =$ \textbf{if} $a$ \textbf{in} $x$ \textbf{then} $T()$ \textbf{else} $T(a)$
	\item $(x \lor T(a)) - y = (x - y) \lor (T(a) - y)$
	\item size$(T()) = 0$
	\item size$(x \lor T(a)) =$ \textbf{if} $a$ \textbf{in} $x$ \textbf{then} size$(x)$ \textbf{else} succ$(\text{size}(x))$

	若定义域类型 $T_0$ 有序，则适用以下公理：

	\item min$(T(a)) = T(a)$
	\item $x \neq T() \supset$ min$(x \lor T(a)) =$ \textbf{if} $a <$ min$(x)$ \textbf{then} $a$ \textbf{else} min$(x)$

	\noindent
	注意：$min(T())$ 未定义。

	\item $x \text{ \textbf{down} } 0 = x \text{ \textbf{up} } 0 = x$
	\item $x$ \textbf{down} succ$(n) = (x$ \textbf{down} $n)$ \textbf{down} 1
	\item $T()$ \textbf{down} $1 = T()$
	\item $(x \lor T(a))$ \textbf{down} $1 = (x$ \textbf{down} $1\ \lor$
	
	\quad \textbf{if} $a \neq T_0$.min \textbf{then} $T($pred$(a))$ \textbf{else} $T()$
	
	\item[(27)\textendash{}(28)] \textbf{up} 类似地定义，把 pred 换成 succ，把 min 换成 max。
	\setcounter{enumi}{28}

	\item $b < a \supset$ range$(a, b) = T()$
	\item $a \leqslant b \supset$ range$(a, b) = T(a)$
	\item $a < b \supset$ range$(a, b) =$ range$(a$, pred$(b)) \lor T(b)$
	\item $j$ \textbf{in} $\{i:D \mid B(i)\} \equiv B(j)$

	\noindent
	简写：
	
	\item $T(a_1, a_2, \dots, a_n)$ 代表 $T(a_1) \lor T(a_2) \lor \dots \lor T(a_n)$
	\item $\{a_1, a_2, \dots, a_n\}$ 代表 $T(a_1, a_2, \dots, a_n)$
	\item $x$: $\land\ y$ 代表 $x \coloneq x \land y$
	\item $x$: $\lor\ y$ 代表 $x \coloneq x \lor y$
	\item $a$ \textbf{from} $x$ 代表 $a \coloneq$ one of $(x)$; $x: - \{a\}$
	\item \textbf{若} $x = \{a_1, a_2, \dots, a_n\}$，\textbf{则}
	
	\qquad\quad \textbf{for} $a$ \textbf{in} $x$ \textbf{do} $S$ 代表
	
	\qquad\quad \quad $S_{a_1}^a; S_{a_2}^a ; \dots; S_{a_n}^a$
	
	\noindent
	其中，如果基类型有序，$a_i$ 按自然顺序排列，否则任意排列；它们不含重复值。
	
	\item $x: -\ y$ 代表 $x \coloneq x - y$
\end{enumerate}

\noindent
定理：
\begin{eqnarray*}
	x = y &\equiv& \forall a:T_0(a \text{ \textbf{in} } x \equiv a \text{ \textbf{in} } y)\\
	&&a \text{ \textbf{in} } (x \lor y) \equiv (a \text{ \textbf{in} } x \lor a \text{ \textbf{in} } y) \\
	&&a \text{ \textbf{in} } (x \land y) \equiv (a \text{ \textbf{in} } x\ \&\ a \text{ \textbf{in} } y) \\
	&&a \text{ \textbf{in} } (x - y) \equiv (a \text{ \textbf{in} } x\ \&\ \neg a \text{ \textbf{in} } y) \\
\end{eqnarray*}

\subsection{序列}

序列定义的一般形式是：

\quad \textbf{type} $T =$ \textbf{sequence} $D$;

\begin{enumerate}[wide, nosep, label=(\arabic*)]
	\item \label{en:sequence-1}
	$T()$ 是 $T$ 类型的值
	
	\item \label{en:sequence-2}
	\textbf{若} $x$ 是 $T$ 类型的值，\textbf{且} $d$ 是 $D$ 类型的值
	
	\qquad\quad \textbf{则} $x\wideparen{\ \ }T(d)$ 是 $T$ 类型的值

	\item $T$ 的元素仅为 \ref{en:sequence-1} 和 \ref{en:sequence-2} 所规定的那些值

	\item $(x\wideparen{\ \ }T(d))$.last $= d$

	\item initial$(x\wideparen{\ \ }T(d)) = x$
	\item $x\wideparen{\ \ }(y\wideparen{\ \ }z) = (x\wideparen{\ \ }y)\wideparen{\ \ }z$
	\item $T(d)$.first $= d$
	\item $x \neq T() \supset (x\wideparen{\ \ }T(d))$.first $= x$.first
	\item final$(T(d)) = T()$
	\item $x \neq T() \supset$ final$(x\wideparen{\ \ }T(d)) =$ final$(x)\wideparen{\ \ }T(d)$

	\noindent
	\textit{注意}：last、initial、first 和 final 对于 $T()$ 未定义。

	\item $T()$ \textbf{ends} $y$
	\item $x\wideparen{\ \ }T(d)$ \textbf{ends} $y \equiv y \neq T()\ \&\ y$.last $= d\ \&\ x$ \textbf{ends} initial$(y)$
	\item $x$ \textbf{begins} $T() \equiv x = T()$
	\item $x$ \textbf{begins} $y\wideparen{\ \ }T(d) \equiv x = y\wideparen{\ \ }T(d) \lor x$ \textbf{begins} $y$
	\item length$(T()) = 0$
	\item length$(x\wideparen{\ \ }T(d)) = \text{succ}(\text{length}(x))$

	\noindent
	对于有序的序列类型，有：

	\item $T() \leqslant y$
	\item $x \leqslant T() \supset x = T()$
	\item $x, y \neq T() \supset (x \leqslant y \equiv x.\text{first} < y.\text{first} \lor (x.\text{first} = y.\text{first}$
	
	\tabto*{12.6em}\& $\text{final}(x) \leqslant \text{final}(y)))$

	\noindent	
	简写：

	\item  $x\wideparen{\text{:} \ \ }T(d)$ \tabto*{11em} 表示 $x \coloneq x\wideparen{\ \ }T(d)$
	\item $d$ \textbf{from} $x$\tabto*{11em} 表示 $d\coloneq x.\text{first};\ x\coloneq \text{final}(x)$
	\item $d$ \textbf{back from} $x$\tabto*{11em} 表示 $d\coloneq x.\text{last};\ x\coloneq \text{initial}(x)$
	\item \textbf{from} $x$\tabto*{11em} 表示 $x\coloneq \text{final}(x)$
	\item \textbf{back from} $x$\tabto*{11em} 表示 $x\coloneq \text{initial}(x)$
	\item $T(d_1, d_2, \dots, d_n)$ 代表
	
	\qquad\quad $(T()\wideparen{\ \ }T(d_1)\wideparen{\ \ }T(d_2)\wideparen{\ \ }\dots\wideparen{\ \ }T(d_n)$
	
	\item $[d_1, d_2, \dots, d_n]$ 代表 $T(d_1, d_2, \dots, d_n)$
	\item \textbf{若} $x = [d_1, d_2, \dots, d_n]$，\textbf{则}
	
	\qquad \quad \textbf{for} $d$ \textbf{in} $x$ \textbf{do} $S$ 代表

	\qquad \qquad \qquad $S_{d_1}^d; S_{d_2}^d; \,.\,.\,; S_{d_n}^d$。


	\qquad \quad \textbf{for} $d$ \textbf{in} $x$ \textbf{take} $E$ 代表
	
	\qquad \qquad \qquad $[E_{d_1}^d, E_{d_2}^d, \dots, E_{d_n}^d]$

\end{enumerate}

\noindent
定理：

	\quad $x = y \equiv (x = y = T() \lor x.\text{first} = y.\text{first}\ \&\ x.\text{final} = y.\text{final})$

	\tabto*{4.9em} $\equiv (x = y = T() \lor x.\text{last} = y.\text{last}\ \&\ x.\text{initial} = y.\text{initial})$
	\bigskip

\noindent
\textbf{参考文献}
\addcontentsline{toc}{section}{参考文献}
\medskip
\nopagebreak

\noindent
以下著作为本章提供了启发和指引，也推荐作为进一步阅读材料。

我还深深感谢 N. Wirth 教授，与他的许多讨论和建议都很有收获；他也愿意在 PASCAL 的设计与实现中检验本文的若干思想。同样感谢 E. W. Dijkstra 教授源源不断的启发。
\medskip

\noindent
Dijkstra, E. W. (1972). Notes on Structured Programming. ``Structured Programming''. pp. 1\textendash{}82. Academic Press, London.
\medskip

\noindent
Knuth, D. E. (1968). ``The Art of Computer Programming'' Vol. 1, Chapter 2. Addison-Wesley, Reading, Mass.
\medskip

\noindent
McCarthy, J. (1963). ``A Basis for a Mathemetical Theory of Computation in Computer Programming and Formal Systems'' (eds. Braffort, P. \& Hirschberg D.). North-Holland Publishing Company, Amsterdam.
\medskip

\noindent
Mealy, G. H. (1967). Another Look at Data. \textit{A.F.I.P.S. Fall Joint Computer Conference Proceedings}. \textbf{31}, pp. 525\textendash{}534.
\medskip

\noindent
Wirth, N. (1970). Programming and Programming Languages. Contribution to Conference of European Chapter of \textit{A.C.M}, Bonn.
\medskip

\noindent
Wirth, N. (1971). Program Development by Stepwise Refinement. \textit{Comm. A.C.M.} \textbf{14}, 4, pp. 221\textendash{}227.
\medskip

\noindent
Wirth, N. (1971). The Programming Language PASCAL. \textit{Acta Informatica}, \textbf{1}, 1, pp. 35\textendash{}63.
